%
% argument.tex -- Argumentpfad
%
% (c) 2021 Prof Dr Andreas Müller, OST Ostschweizer Fachhochschule
%
\documentclass[tikz]{standalone}
\usepackage{amsmath}
\usepackage{times}
\usepackage{txfonts}
\usepackage{pgfplots}
\usepackage{csvsimple}
\usetikzlibrary{arrows,intersections,math,calc}
\definecolor{darkred}{rgb}{0.8,0,0}
\definecolor{darkgreen}{rgb}{0,0.6,0}
\begin{document}
\def\skala{1}
%
% argumentpfad.tex -- generated by argument.m, don't edit
%
% (c) 2026 Prof Dr Andreas Müller
%
\def\nullstelleeins{(1.0000,0.0000)}
\def\nullstellezwei{(-0.5000,0.5000)}
\def\nullstelledrei{(-0.2000,-0.7000)}
\def\poleins{(0.0000,0.8500)}
\def\polzwei{(1.3000,-1.5000)}
\def\r{1.3000}
\def\bildkurveA{
	(0.3000,0.0000)
	-- (0.2990,0.0514)
	-- (0.2959,0.1026)
	-- (0.2909,0.1538)
	-- (0.2838,0.2046)
	-- (0.2747,0.2552)
	-- (0.2636,0.3054)
	-- (0.2506,0.3550)
	-- (0.2356,0.4042)
	-- (0.2186,0.4527)
	-- (0.1998,0.5005)
	-- (0.1791,0.5475)
	-- (0.1565,0.5936)
	-- (0.1322,0.6389)
	-- (0.1061,0.6831)
	-- (0.0782,0.7263)
	-- (0.0487,0.7683)
	-- (0.0175,0.8091)
	-- (-0.0153,0.8487)
	-- (-0.0495,0.8869)
	-- (-0.0853,0.9238)
	-- (-0.1225,0.9592)
	-- (-0.1611,0.9931)
	-- (-0.2010,1.0255)
	-- (-0.2422,1.0563)
	-- (-0.2845,1.0854)
	-- (-0.3279,1.1128)
	-- (-0.3724,1.1385)
	-- (-0.4179,1.1624)
	-- (-0.4642,1.1845)
	-- (-0.5115,1.2047)
	-- (-0.5594,1.2231)
	-- (-0.6081,1.2395)
	-- (-0.6574,1.2540)
	-- (-0.7072,1.2666)
	-- (-0.7575,1.2772)
	-- (-0.8081,1.2858)
	-- (-0.8591,1.2923)
	-- (-0.9102,1.2969)
	-- (-0.9615,1.2994)
	-- (-1.0129,1.2999)
	-- (-1.0642,1.2984)
	-- (-1.1155,1.2949)
	-- (-1.1666,1.2893)
	-- (-1.2174,1.2817)
	-- (-1.2678,1.2721)
	-- (-1.3179,1.2605)
	-- (-1.3674,1.2470)
	-- (-1.4164,1.2315)
	-- (-1.4647,1.2141)
	-- (-1.5123,1.1948)
	-- (-1.5591,1.1736)
	-- (-1.6051,1.1506)
	-- (-1.6501,1.1258)
	-- (-1.6940,1.0992)
	-- (-1.7369,1.0710)
	-- (-1.7787,1.0410)
	-- (-1.8192,1.0094)
	-- (-1.8584,0.9763)
	-- (-1.8963,0.9416)
	-- (-1.9328,0.9055)
	-- (-1.9679,0.8679)
	-- (-2.0014,0.8290)
	-- (-2.0334,0.7888)
	-- (-2.0637,0.7473)
	-- (-2.0924,0.7047)
	-- (-2.1194,0.6610)
	-- (-2.1446,0.6163)
	-- (-2.1681,0.5706)
	-- (-2.1897,0.5240)
	-- (-2.2095,0.4766)
	-- (-2.2274,0.4284)
	-- (-2.2433,0.3796)
	-- (-2.2574,0.3302)
	-- (-2.2694,0.2802)
	-- (-2.2795,0.2299)
	-- (-2.2876,0.1791)
	-- (-2.2937,0.1281)
	-- (-2.2977,0.0769)
	-- (-2.2997,0.0256)
	-- (-2.2997,-0.0258)
	-- (-2.2977,-0.0771)
	-- (-2.2937,-0.1283)
	-- (-2.2876,-0.1793)
	-- (-2.2795,-0.2301)
	-- (-2.2694,-0.2804)
	-- (-2.2573,-0.3304)
	-- (-2.2433,-0.3798)
	-- (-2.2273,-0.4286)
	-- (-2.2094,-0.4768)
	-- (-2.1896,-0.5242)
	-- (-2.1680,-0.5708)
	-- (-2.1445,-0.6165)
	-- (-2.1193,-0.6612)
	-- (-2.0923,-0.7049)
	-- (-2.0636,-0.7475)
	-- (-2.0332,-0.7889)
	-- (-2.0013,-0.8291)
	-- (-1.9677,-0.8681)
	-- (-1.9327,-0.9056)
	-- (-1.8962,-0.9418)
	-- (-1.8583,-0.9764)
	-- (-1.8190,-1.0096)
	-- (-1.7785,-1.0411)
	-- (-1.7367,-1.0711)
	-- (-1.6938,-1.0994)
	-- (-1.6499,-1.1259)
	-- (-1.6049,-1.1507)
	-- (-1.5590,-1.1737)
	-- (-1.5121,-1.1949)
	-- (-1.4645,-1.2142)
	-- (-1.4162,-1.2316)
	-- (-1.3672,-1.2471)
	-- (-1.3177,-1.2606)
	-- (-1.2676,-1.2722)
	-- (-1.2172,-1.2817)
	-- (-1.1663,-1.2893)
	-- (-1.1153,-1.2949)
	-- (-1.0640,-1.2984)
	-- (-1.0127,-1.2999)
	-- (-0.9613,-1.2994)
	-- (-0.9100,-1.2969)
	-- (-0.8588,-1.2923)
	-- (-0.8079,-1.2857)
	-- (-0.7572,-1.2771)
	-- (-0.7070,-1.2665)
	-- (-0.6572,-1.2540)
	-- (-0.6079,-1.2395)
	-- (-0.5592,-1.2230)
	-- (-0.5113,-1.2046)
	-- (-0.4641,-1.1844)
	-- (-0.4177,-1.1623)
	-- (-0.3722,-1.1384)
	-- (-0.3277,-1.1127)
	-- (-0.2843,-1.0852)
	-- (-0.2420,-1.0561)
	-- (-0.2008,-1.0254)
	-- (-0.1610,-0.9930)
	-- (-0.1224,-0.9591)
	-- (-0.0852,-0.9236)
	-- (-0.0494,-0.8868)
	-- (-0.0151,-0.8485)
	-- (0.0176,-0.8090)
	-- (0.0488,-0.7681)
	-- (0.0783,-0.7261)
	-- (0.1062,-0.6829)
	-- (0.1323,-0.6387)
	-- (0.1566,-0.5934)
	-- (0.1792,-0.5473)
	-- (0.1999,-0.5003)
	-- (0.2187,-0.4525)
	-- (0.2356,-0.4040)
	-- (0.2506,-0.3548)
	-- (0.2637,-0.3051)
	-- (0.2747,-0.2550)
	-- (0.2838,-0.2044)
	-- (0.2909,-0.1535)
	-- (0.2960,-0.1024)
	-- (0.2990,-0.0511)
	-- (0.3000,0.0002)
	-- (0.2990,0.0516)
	-- (0.2959,0.1029)
	-- (0.2908,0.1540)
	-- (0.2838,0.2049)
	-- (0.2747,0.2554)
	-- (0.2636,0.3056)
	-- (0.2505,0.3553)
	-- (0.2355,0.4044)
	-- (0.2186,0.4529)
	-- (0.1997,0.5007)
	-- (0.1790,0.5477)
	-- (0.1564,0.5938)
	-- (0.1321,0.6390)
	-- (0.1060,0.6833)
	-- (0.0781,0.7264)
	-- (0.0486,0.7685)
	-- (0.0174,0.8093)
	-- (-0.0154,0.8489)
	-- (-0.0497,0.8871)
	-- (-0.0855,0.9239)
	-- (-0.1227,0.9593)
	-- (-0.1613,0.9933)
	-- (-0.2012,1.0256)
	-- (-0.2423,1.0564)
	-- (-0.2847,1.0855)
	-- (-0.3281,1.1129)
	-- (-0.3726,1.1386)
	-- (-0.4181,1.1625)
	-- (-0.4644,1.1846)
	-- (-0.5117,1.2048)
	-- (-0.5596,1.2231)
	-- (-0.6083,1.2396)
	-- (-0.6576,1.2541)
	-- (-0.7074,1.2666)
	-- (-0.7577,1.2772)
	-- (-0.8083,1.2858)
	-- (-0.8593,1.2924)
	-- (-0.9104,1.2969)
	-- (-0.9617,1.2994)
	-- (-1.0131,1.2999)
	-- (-1.0645,1.2984)
	-- (-1.1157,1.2948)
	-- (-1.1668,1.2893)
	-- (-1.2176,1.2817)
	-- (-1.2680,1.2721)
	-- (-1.3181,1.2605)
	-- (-1.3676,1.2469)
	-- (-1.4166,1.2314)
	-- (-1.4649,1.2140)
	-- (-1.5125,1.1947)
	-- (-1.5593,1.1735)
	-- (-1.6053,1.1505)
	-- (-1.6502,1.1257)
	-- (-1.6942,1.0991)
	-- (-1.7371,1.0708)
	-- (-1.7788,1.0409)
	-- (-1.8193,1.0093)
	-- (-1.8586,0.9761)
	-- (-1.8965,0.9415)
	-- (-1.9330,0.9053)
	-- (-1.9680,0.8677)
	-- (-2.0015,0.8288)
	-- (-2.0335,0.7886)
	-- (-2.0638,0.7472)
	-- (-2.0925,0.7045)
	-- (-2.1195,0.6608)
	-- (-2.1447,0.6161)
	-- (-2.1682,0.5704)
	-- (-2.1898,0.5238)
	-- (-2.2096,0.4764)
	-- (-2.2275,0.4282)
	-- (-2.2434,0.3794)
	-- (-2.2574,0.3300)
	-- (-2.2695,0.2800)
	-- (-2.2796,0.2296)
	-- (-2.2876,0.1789)
	-- (-2.2937,0.1279)
	-- (-2.2977,0.0767)
	-- (-2.2998,0.0254)
	-- (-2.2997,-0.0260)
	-- (-2.2977,-0.0773)
	-- (-2.2936,-0.1285)
	-- (-2.2875,-0.1796)
	-- (-2.2794,-0.2303)
	-- (-2.2693,-0.2807)
	-- (-2.2573,-0.3306)
	-- (-2.2432,-0.3800)
	-- (-2.2272,-0.4288)
	-- (-2.2093,-0.4770)
	-- (-2.1896,-0.5244)
	-- (-2.1679,-0.5710)
	-- (-2.1444,-0.6167)
	-- (-2.1192,-0.6614)
	-- (-2.0922,-0.7051)
	-- (-2.0635,-0.7477)
	-- (-2.0331,-0.7891)
	-- (-2.0011,-0.8293)
	-- (-1.9676,-0.8682)
	-- (-1.9325,-0.9058)
	-- (-1.8960,-0.9419)
	-- (-1.8581,-0.9766)
	-- (-1.8188,-1.0097)
	-- (-1.7783,-1.0413)
	-- (-1.7366,-1.0712)
	-- (-1.6937,-1.0995)
	-- (-1.6497,-1.1260)
	-- (-1.6047,-1.1508)
	-- (-1.5588,-1.1738)
	-- (-1.5119,-1.1950)
	-- (-1.4643,-1.2142)
	-- (-1.4160,-1.2316)
	-- (-1.3670,-1.2471)
	-- (-1.3175,-1.2606)
	-- (-1.2674,-1.2722)
	-- (-1.2169,-1.2818)
	-- (-1.1661,-1.2893)
	-- (-1.1151,-1.2949)
	-- (-1.0638,-1.2984)
	-- (-1.0125,-1.2999)
	-- (-0.9611,-1.2994)
	-- (-0.9098,-1.2969)
	-- (-0.8586,-1.2923)
	-- (-0.8077,-1.2857)
	-- (-0.7570,-1.2771)
	-- (-0.7068,-1.2665)
	-- (-0.6570,-1.2539)
	-- (-0.6077,-1.2394)
	-- (-0.5590,-1.2229)
	-- (-0.5111,-1.2045)
	-- (-0.4639,-1.1843)
	-- (-0.4175,-1.1622)
	-- (-0.3720,-1.1383)
	-- (-0.3275,-1.1126)
	-- (-0.2841,-1.0851)
	-- (-0.2418,-1.0560)
	-- (-0.2007,-1.0252)
	-- (-0.1608,-0.9928)
	-- (-0.1222,-0.9589)
	-- (-0.0850,-0.9235)
	-- (-0.0493,-0.8866)
	-- (-0.0150,-0.8484)
	-- (0.0178,-0.8088)
	-- (0.0489,-0.7679)
	-- (0.0785,-0.7259)
	-- (0.1063,-0.6827)
	-- (0.1324,-0.6385)
	-- (0.1567,-0.5932)
	-- (0.1793,-0.5471)
	-- (0.2000,-0.5001)
	-- (0.2188,-0.4523)
	-- (0.2357,-0.4038)
	-- (0.2507,-0.3546)
	-- (0.2637,-0.3049)
	-- (0.2748,-0.2548)
	-- (0.2839,-0.2042)
	-- (0.2909,-0.1533)
	-- (0.2960,-0.1022)
	-- (0.2990,-0.0509)
	-- (0.3000,0.0004)
	-- (0.2990,0.0518)
	-- (0.2959,0.1031)
	-- (0.2908,0.1542)
	-- (0.2837,0.2051)
	-- (0.2746,0.2556)
	-- (0.2635,0.3058)
	-- (0.2505,0.3555)
	-- (0.2354,0.4046)
	-- (0.2185,0.4531)
	-- (0.1996,0.5009)
	-- (0.1789,0.5479)
	-- (0.1563,0.5940)
	-- (0.1320,0.6392)
	-- (0.1058,0.6835)
	-- (0.0780,0.7266)
	-- (0.0484,0.7686)
	-- (0.0172,0.8095)
	-- (-0.0155,0.8490)
	-- (-0.0498,0.8872)
	-- (-0.0856,0.9241)
	-- (-0.1229,0.9595)
	-- (-0.1615,0.9934)
	-- (-0.2014,1.0258)
	-- (-0.2425,1.0565)
	-- (-0.2848,1.0856)
	-- (-0.3283,1.1130)
	-- (-0.3728,1.1387)
	-- (-0.4183,1.1626)
	-- (-0.4646,1.1846)
	-- (-0.5119,1.2049)
	-- (-0.5598,1.2232)
	-- (-0.6085,1.2397)
	-- (-0.6578,1.2542)
	-- (-0.7076,1.2667)
	-- (-0.7579,1.2773)
	-- (-0.8085,1.2858)
	-- (-0.8595,1.2924)
	-- (-0.9107,1.2969)
	-- (-0.9620,1.2994)
	-- (-1.0133,1.2999)
	-- (-1.0647,1.2984)
	-- (-1.1159,1.2948)
	-- (-1.1670,1.2892)
	-- (-1.2178,1.2816)
	-- (-1.2683,1.2720)
	-- (-1.3183,1.2604)
	-- (-1.3678,1.2469)
	-- (-1.4168,1.2314)
	-- (-1.4651,1.2139)
	-- (-1.5127,1.1946)
	-- (-1.5595,1.1734)
	-- (-1.6055,1.1504)
	-- (-1.6504,1.1256)
	-- (-1.6944,1.0990)
	-- (-1.7373,1.0707)
	-- (-1.7790,1.0407)
	-- (-1.8195,1.0092)
	-- (-1.8587,0.9760)
	-- (-1.8966,0.9413)
	-- (-1.9331,0.9051)
	-- (-1.9681,0.8676)
	-- (-2.0017,0.8286)
	-- (-2.0336,0.7884)
	-- (-2.0640,0.7470)
	-- (-2.0926,0.7044)
	-- (-2.1196,0.6606)
	-- (-2.1448,0.6159)
	-- (-2.1683,0.5702)
	-- (-2.1899,0.5236)
	-- (-2.2097,0.4762)
	-- (-2.2275,0.4280)
	-- (-2.2435,0.3792)
	-- (-2.2575,0.3297)
	-- (-2.2695,0.2798)
	-- (-2.2796,0.2294)
	-- (-2.2877,0.1787)
	-- (-2.2937,0.1277)
	-- (-2.2977,0.0765)
	-- (-2.2998,0.0251)
	-- (-2.2997,-0.0262)
	-- (-2.2977,-0.0776)
	-- (-2.2936,-0.1288)
	-- (-2.2875,-0.1798)
	-- (-2.2794,-0.2305)
	-- (-2.2693,-0.2809)
	-- (-2.2572,-0.3308)
	-- (-2.2432,-0.3802)
	-- (-2.2272,-0.4290)
	-- (-2.2093,-0.4772)
	-- (-2.1895,-0.5246)
	-- (-2.1678,-0.5712)
	-- (-2.1443,-0.6168)
	-- (-2.1191,-0.6616)
	-- (-2.0921,-0.7053)
	-- (-2.0633,-0.7479)
	-- (-2.0330,-0.7893)
	-- (-2.0010,-0.8295)
	-- (-1.9674,-0.8684)
	-- (-1.9324,-0.9059)
	-- (-1.8958,-0.9421)
	-- (-1.8579,-0.9767)
	-- (-1.8187,-1.0098)
	-- (-1.7781,-1.0414)
	-- (-1.7364,-1.0713)
	-- (-1.6935,-1.0996)
	-- (-1.6495,-1.1261)
	-- (-1.6045,-1.1509)
	-- (-1.5586,-1.1739)
	-- (-1.5117,-1.1950)
	-- (-1.4641,-1.2143)
	-- (-1.4158,-1.2317)
	-- (-1.3668,-1.2472)
	-- (-1.3173,-1.2607)
	-- (-1.2672,-1.2722)
	-- (-1.2167,-1.2818)
	-- (-1.1659,-1.2894)
	-- (-1.1148,-1.2949)
	-- (-1.0636,-1.2984)
	-- (-1.0122,-1.2999)
	-- (-0.9609,-1.2994)
	-- (-0.9096,-1.2969)
	-- (-0.8584,-1.2923)
	-- (-0.8075,-1.2857)
	-- (-0.7568,-1.2771)
	-- (-0.7066,-1.2664)
	-- (-0.6568,-1.2539)
	-- (-0.6075,-1.2393)
	-- (-0.5588,-1.2229)
	-- (-0.5109,-1.2045)
	-- (-0.4637,-1.1842)
	-- (-0.4173,-1.1621)
	-- (-0.3718,-1.1382)
	-- (-0.3274,-1.1125)
	-- (-0.2839,-1.0850)
	-- (-0.2416,-1.0559)
	-- (-0.2005,-1.0251)
	-- (-0.1606,-0.9927)
	-- (-0.1221,-0.9588)
	-- (-0.0849,-0.9233)
	-- (-0.0491,-0.8865)
	-- (-0.0148,-0.8482)
	-- (0.0179,-0.8086)
	-- (0.0491,-0.7678)
	-- (0.0786,-0.7257)
	-- (0.1064,-0.6825)
	-- (0.1325,-0.6383)
	-- (0.1568,-0.5931)
	-- (0.1794,-0.5469)
	-- (0.2001,-0.4999)
	-- (0.2189,-0.4521)
	-- (0.2358,-0.4036)
	-- (0.2508,-0.3544)
	-- (0.2638,-0.3047)
	-- (0.2748,-0.2546)
	-- (0.2839,-0.2040)
	-- (0.2910,-0.1531)
	-- (0.2960,-0.1020)
	-- (0.2990,-0.0507)
	-- (0.3000,0.0006)
	-- (0.2990,0.0520)
	-- (0.2959,0.1033)
	-- (0.2908,0.1544)
	-- (0.2837,0.2053)
	-- (0.2746,0.2558)
	-- (0.2635,0.3060)
	-- (0.2504,0.3557)
	-- (0.2354,0.4048)
	-- (0.2184,0.4533)
	-- (0.1996,0.5011)
	-- (0.1788,0.5481)
	-- (0.1563,0.5942)
	-- (0.1319,0.6394)
	-- (0.1057,0.6836)
	-- (0.0779,0.7268)
	-- (0.0483,0.7688)
	-- (0.0171,0.8096)
	-- (-0.0157,0.8492)
	-- (-0.0500,0.8874)
	-- (-0.0858,0.9242)
	-- (-0.1230,0.9596)
	-- (-0.1616,0.9935)
	-- (-0.2015,1.0259)
	-- (-0.2427,1.0566)
	-- (-0.2850,1.0857)
	-- (-0.3285,1.1131)
	-- (-0.3730,1.1388)
	-- (-0.4185,1.1627)
	-- (-0.4648,1.1847)
	-- (-0.5121,1.2050)
	-- (-0.5600,1.2233)
	-- (-0.6087,1.2397)
	-- (-0.6580,1.2542)
	-- (-0.7078,1.2667)
	-- (-0.7581,1.2773)
	-- (-0.8087,1.2859)
	-- (-0.8597,1.2924)
	-- (-0.9109,1.2969)
	-- (-0.9622,1.2994)
	-- (-1.0135,1.2999)
	-- (-1.0649,1.2984)
	-- (-1.1161,1.2948)
	-- (-1.1672,1.2892)
	-- (-1.2180,1.2816)
	-- (-1.2685,1.2720)
	-- (-1.3185,1.2604)
	-- (-1.3681,1.2468)
	-- (-1.4170,1.2313)
	-- (-1.4653,1.2139)
	-- (-1.5129,1.1945)
	-- (-1.5597,1.1733)
	-- (-1.6057,1.1503)
	-- (-1.6506,1.1255)
	-- (-1.6946,1.0989)
	-- (-1.7375,1.0706)
	-- (-1.7792,1.0406)
	-- (-1.8197,1.0090)
	-- (-1.8589,0.9759)
	-- (-1.8968,0.9412)
	-- (-1.9333,0.9050)
	-- (-1.9683,0.8674)
	-- (-2.0018,0.8285)
	-- (-2.0338,0.7883)
	-- (-2.0641,0.7468)
	-- (-2.0928,0.7042)
	-- (-2.1197,0.6605)
	-- (-2.1450,0.6157)
	-- (-2.1684,0.5700)
	-- (-2.1900,0.5234)
	-- (-2.2097,0.4760)
	-- (-2.2276,0.4278)
	-- (-2.2435,0.3790)
	-- (-2.2575,0.3295)
	-- (-2.2696,0.2796)
	-- (-2.2796,0.2292)
	-- (-2.2877,0.1785)
	-- (-2.2937,0.1275)
	-- (-2.2978,0.0763)
	-- (-2.2998,0.0249)
	-- (-2.2997,-0.0264)
	-- (-2.2977,-0.0778)
	-- (-2.2936,-0.1290)
	-- (-2.2875,-0.1800)
	-- (-2.2794,-0.2307)
	-- (-2.2693,-0.2811)
	-- (-2.2572,-0.3310)
	-- (-2.2431,-0.3804)
	-- (-2.2271,-0.4292)
	-- (-2.2092,-0.4774)
	-- (-2.1894,-0.5248)
	-- (-2.1677,-0.5713)
	-- (-2.1442,-0.6170)
	-- (-2.1190,-0.6618)
	-- (-2.0919,-0.7054)
	-- (-2.0632,-0.7480)
	-- (-2.0328,-0.7895)
	-- (-2.0008,-0.8296)
	-- (-1.9673,-0.8685)
	-- (-1.9322,-0.9061)
	-- (-1.8957,-0.9422)
	-- (-1.8578,-0.9769)
	-- (-1.8185,-1.0100)
	-- (-1.7780,-1.0415)
	-- (-1.7362,-1.0714)
	-- (-1.6933,-1.0997)
	-- (-1.6493,-1.1262)
	-- (-1.6043,-1.1510)
	-- (-1.5584,-1.1740)
	-- (-1.5115,-1.1951)
	-- (-1.4639,-1.2144)
	-- (-1.4156,-1.2318)
	-- (-1.3666,-1.2472)
	-- (-1.3170,-1.2607)
	-- (-1.2670,-1.2723)
	-- (-1.2165,-1.2818)
	-- (-1.1657,-1.2894)
	-- (-1.1146,-1.2949)
	-- (-1.0634,-1.2985)
	-- (-1.0120,-1.2999)
	-- (-0.9607,-1.2994)
	-- (-0.9094,-1.2968)
	-- (-0.8582,-1.2922)
	-- (-0.8073,-1.2856)
	-- (-0.7566,-1.2770)
	-- (-0.7063,-1.2664)
	-- (-0.6565,-1.2538)
	-- (-0.6073,-1.2393)
	-- (-0.5586,-1.2228)
	-- (-0.5107,-1.2044)
	-- (-0.4635,-1.1841)
	-- (-0.4171,-1.1620)
	-- (-0.3716,-1.1381)
	-- (-0.3272,-1.1123)
	-- (-0.2838,-1.0849)
	-- (-0.2415,-1.0557)
	-- (-0.2003,-1.0250)
	-- (-0.1605,-0.9926)
	-- (-0.1219,-0.9586)
	-- (-0.0847,-0.9232)
	-- (-0.0490,-0.8863)
	-- (-0.0147,-0.8480)
	-- (0.0180,-0.8084)
	-- (0.0492,-0.7676)
	-- (0.0787,-0.7255)
	-- (0.1065,-0.6824)
	-- (0.1326,-0.6381)
	-- (0.1569,-0.5929)
	-- (0.1795,-0.5467)
	-- (0.2001,-0.4997)
	-- (0.2189,-0.4519)
	-- (0.2358,-0.4034)
	-- (0.2508,-0.3542)
	-- (0.2638,-0.3045)
	-- (0.2749,-0.2543)
	-- (0.2839,-0.2038)
	-- (0.2910,-0.1529)
	-- (0.2960,-0.1018)
	-- (0.2990,-0.0505)
	-- (0.3000,0.0009)
	-- (0.2990,0.0522)
	-- (0.2959,0.1035)
	-- (0.2908,0.1546)
	-- (0.2837,0.2055)
	-- (0.2745,0.2560)
	-- (0.2634,0.3062)
	-- (0.2503,0.3559)
	-- (0.2353,0.4050)
	-- (0.2183,0.4535)
	-- (0.1995,0.5013)
	-- (0.1787,0.5483)
	-- (0.1562,0.5944)
	-- (0.1318,0.6396)
	-- (0.1056,0.6838)
	-- (0.0777,0.7270)
	-- (0.0482,0.7690)
	-- (0.0170,0.8098)
	-- (-0.0158,0.8493)
	-- (-0.0501,0.8876)
	-- (-0.0859,0.9244)
	-- (-0.1232,0.9598)
	-- (-0.1618,0.9937)
	-- (-0.2017,1.0260)
	-- (-0.2429,1.0568)
	-- (-0.2852,1.0858)
	-- (-0.3287,1.1132)
	-- (-0.3732,1.1389)
	-- (-0.4186,1.1628)
	-- (-0.4650,1.1848)
	-- (-0.5123,1.2050)
	-- (-0.5603,1.2234)
	-- (-0.6089,1.2398)
	-- (-0.6582,1.2543)
	-- (-0.7080,1.2668)
	-- (-0.7583,1.2773)
	-- (-0.8090,1.2859)
	-- (-0.8599,1.2924)
	-- (-0.9111,1.2970)
	-- (-0.9624,1.2995)
	-- (-1.0138,1.2999)
	-- (-1.0651,1.2984)
	-- (-1.1164,1.2948)
	-- (-1.1674,1.2892)
	-- (-1.2182,1.2816)
	-- (-1.2687,1.2719)
	-- (-1.3187,1.2603)
	-- (-1.3683,1.2467)
	-- (-1.4172,1.2312)
	-- (-1.4655,1.2138)
	-- (-1.5131,1.1944)
	-- (-1.5599,1.1732)
	-- (-1.6058,1.1502)
	-- (-1.6508,1.1254)
	-- (-1.6948,1.0988)
	-- (-1.7376,1.0705)
	-- (-1.7793,1.0405)
	-- (-1.8198,1.0089)
	-- (-1.8591,0.9757)
	-- (-1.8969,0.9410)
	-- (-1.9334,0.9048)
	-- (-1.9684,0.8673)
	-- (-2.0019,0.8283)
	-- (-2.0339,0.7881)
	-- (-2.0642,0.7466)
	-- (-2.0929,0.7040)
	-- (-2.1198,0.6603)
	-- (-2.1451,0.6155)
	-- (-2.1685,0.5698)
	-- (-2.1901,0.5232)
	-- (-2.2098,0.4758)
	-- (-2.2277,0.4276)
	-- (-2.2436,0.3788)
	-- (-2.2576,0.3293)
	-- (-2.2696,0.2794)
	-- (-2.2797,0.2290)
	-- (-2.2877,0.1783)
	-- (-2.2938,0.1273)
	-- (-2.2978,0.0760)
	-- (-2.2998,0.0247)
	-- (-2.2997,-0.0267)
	-- (-2.2977,-0.0780)
	-- (-2.2936,-0.1292)
	-- (-2.2875,-0.1802)
	-- (-2.2793,-0.2309)
	-- (-2.2692,-0.2813)
	-- (-2.2571,-0.3312)
	-- (-2.2430,-0.3806)
	-- (-2.2270,-0.4294)
	-- (-2.2091,-0.4776)
	-- (-2.1893,-0.5250)
	-- (-2.1676,-0.5715)
	-- (-2.1441,-0.6172)
	-- (-2.1189,-0.6619)
	-- (-2.0918,-0.7056)
	-- (-2.0631,-0.7482)
	-- (-2.0327,-0.7896)
	-- (-2.0007,-0.8298)
	-- (-1.9671,-0.8687)
	-- (-1.9321,-0.9062)
	-- (-1.8955,-0.9423)
	-- (-1.8576,-0.9770)
	-- (-1.8183,-1.0101)
	-- (-1.7778,-1.0417)
	-- (-1.7360,-1.0716)
	-- (-1.6931,-1.0998)
	-- (-1.6491,-1.1263)
	-- (-1.6041,-1.1511)
	-- (-1.5582,-1.1741)
	-- (-1.5113,-1.1952)
	-- (-1.4637,-1.2145)
	-- (-1.4154,-1.2318)
	-- (-1.3664,-1.2473)
	-- (-1.3168,-1.2608)
	-- (-1.2668,-1.2723)
	-- (-1.2163,-1.2819)
	-- (-1.1655,-1.2894)
	-- (-1.1144,-1.2950)
	-- (-1.0632,-1.2985)
	-- (-1.0118,-1.2999)
	-- (-0.9604,-1.2994)
	-- (-0.9091,-1.2968)
	-- (-0.8580,-1.2922)
	-- (-0.8070,-1.2856)
	-- (-0.7564,-1.2770)
	-- (-0.7061,-1.2664)
	-- (-0.6563,-1.2538)
	-- (-0.6071,-1.2392)
	-- (-0.5584,-1.2227)
	-- (-0.5105,-1.2043)
	-- (-0.4633,-1.1840)
	-- (-0.4169,-1.1619)
	-- (-0.3715,-1.1380)
	-- (-0.3270,-1.1122)
	-- (-0.2836,-1.0848)
	-- (-0.2413,-1.0556)
	-- (-0.2002,-1.0248)
	-- (-0.1603,-0.9924)
	-- (-0.1217,-0.9585)
	-- (-0.0846,-0.9230)
	-- (-0.0488,-0.8861)
	-- (-0.0145,-0.8479)
	-- (0.0182,-0.8083)
	-- (0.0493,-0.7674)
	-- (0.0788,-0.7254)
	-- (0.1066,-0.6822)
	-- (0.1327,-0.6379)
	-- (0.1570,-0.5927)
	-- (0.1796,-0.5465)
	-- (0.2002,-0.4995)
	-- (0.2190,-0.4517)
	-- (0.2359,-0.4031)
	-- (0.2509,-0.3540)
	-- (0.2639,-0.3043)
	-- (0.2749,-0.2541)
	-- (0.2840,-0.2036)
	-- (0.2910,-0.1527)
	-- (0.2960,-0.1016)
	-- (0.2990,-0.0503)
	-- (0.3000,0.0011)
	-- (0.2989,0.0524)
	-- (0.2959,0.1037)
	-- (0.2907,0.1548)
	-- (0.2836,0.2057)
	-- (0.2745,0.2563)
	-- (0.2634,0.3064)
	-- (0.2503,0.3561)
	-- (0.2352,0.4052)
	-- (0.2183,0.4537)
	-- (0.1994,0.5015)
	-- (0.1786,0.5485)
	-- (0.1561,0.5946)
	-- (0.1317,0.6398)
	-- (0.1055,0.6840)
	-- (0.0776,0.7272)
	-- (0.0480,0.7692)
	-- (0.0168,0.8100)
	-- (-0.0160,0.8495)
	-- (-0.0503,0.8877)
	-- (-0.0861,0.9245)
	-- (-0.1233,0.9599)
	-- (-0.1620,0.9938)
	-- (-0.2019,1.0262)
	-- (-0.2430,1.0569)
	-- (-0.2854,1.0860)
	-- (-0.3288,1.1133)
	-- (-0.3733,1.1390)
	-- (-0.4188,1.1629)
	-- (-0.4652,1.1849)
	-- (-0.5125,1.2051)
	-- (-0.5605,1.2234)
	-- (-0.6091,1.2398)
	-- (-0.6584,1.2543)
	-- (-0.7082,1.2668)
	-- (-0.7585,1.2774)
	-- (-0.8092,1.2859)
	-- (-0.8601,1.2925)
	-- (-0.9113,1.2970)
	-- (-0.9626,1.2995)
	-- (-1.0140,1.2999)
	-- (-1.0653,1.2984)
	-- (-1.1166,1.2948)
	-- (-1.1676,1.2891)
	-- (-1.2184,1.2815)
	-- (-1.2689,1.2719)
	-- (-1.3189,1.2603)
	-- (-1.3685,1.2467)
	-- (-1.4174,1.2312)
	-- (-1.4658,1.2137)
	-- (-1.5133,1.1944)
	-- (-1.5601,1.1731)
	-- (-1.6060,1.1501)
	-- (-1.6510,1.1253)
	-- (-1.6949,1.0987)
	-- (-1.7378,1.0703)
	-- (-1.7795,1.0404)
	-- (-1.8200,1.0088)
	-- (-1.8592,0.9756)
	-- (-1.8971,0.9409)
	-- (-1.9336,0.9047)
	-- (-1.9686,0.8671)
	-- (-2.0021,0.8281)
	-- (-2.0340,0.7879)
	-- (-2.0643,0.7464)
	-- (-2.0930,0.7038)
	-- (-2.1200,0.6601)
	-- (-2.1452,0.6153)
	-- (-2.1686,0.5696)
	-- (-2.1902,0.5230)
	-- (-2.2099,0.4756)
	-- (-2.2277,0.4274)
	-- (-2.2437,0.3785)
	-- (-2.2576,0.3291)
	-- (-2.2697,0.2792)
	-- (-2.2797,0.2288)
	-- (-2.2877,0.1781)
	-- (-2.2938,0.1270)
	-- (-2.2978,0.0758)
	-- (-2.2998,0.0245)
	-- (-2.2997,-0.0269)
	-- (-2.2976,-0.0782)
	-- (-2.2935,-0.1294)
	-- (-2.2874,-0.1804)
	-- (-2.2793,-0.2311)
	-- (-2.2692,-0.2815)
	-- (-2.2570,-0.3314)
	-- (-2.2430,-0.3808)
	-- (-2.2270,-0.4296)
	-- (-2.2090,-0.4778)
	-- (-2.1892,-0.5252)
	-- (-2.1675,-0.5717)
	-- (-2.1440,-0.6174)
	-- (-2.1187,-0.6621)
	-- (-2.0917,-0.7058)
	-- (-2.0630,-0.7484)
	-- (-2.0326,-0.7898)
	-- (-2.0006,-0.8300)
	-- (-1.9670,-0.8689)
	-- (-1.9319,-0.9064)
	-- (-1.8954,-0.9425)
	-- (-1.8574,-0.9771)
	-- (-1.8182,-1.0102)
	-- (-1.7776,-1.0418)
	-- (-1.7358,-1.0717)
	-- (-1.6929,-1.0999)
	-- (-1.6489,-1.1264)
	-- (-1.6039,-1.1512)
	-- (-1.5580,-1.1742)
	-- (-1.5112,-1.1953)
	-- (-1.4635,-1.2146)
	-- (-1.4152,-1.2319)
	-- (-1.3662,-1.2474)
	-- (-1.3166,-1.2609)
	-- (-1.2666,-1.2724)
	-- (-1.2161,-1.2819)
	-- (-1.1653,-1.2895)
	-- (-1.1142,-1.2950)
	-- (-1.0629,-1.2985)
	-- (-1.0116,-1.2999)
	-- (-0.9602,-1.2994)
	-- (-0.9089,-1.2968)
	-- (-0.8578,-1.2922)
	-- (-0.8068,-1.2856)
	-- (-0.7562,-1.2769)
	-- (-0.7059,-1.2663)
	-- (-0.6561,-1.2537)
	-- (-0.6069,-1.2391)
	-- (-0.5582,-1.2226)
	-- (-0.5103,-1.2042)
	-- (-0.4631,-1.1839)
	-- (-0.4167,-1.1618)
	-- (-0.3713,-1.1378)
	-- (-0.3268,-1.1121)
	-- (-0.2834,-1.0847)
	-- (-0.2411,-1.0555)
	-- (-0.2000,-1.0247)
	-- (-0.1601,-0.9923)
	-- (-0.1216,-0.9583)
	-- (-0.0844,-0.9229)
	-- (-0.0487,-0.8860)
	-- (-0.0144,-0.8477)
	-- (0.0183,-0.8081)
	-- (0.0494,-0.7672)
	-- (0.0789,-0.7252)
	-- (0.1067,-0.6820)
	-- (0.1328,-0.6377)
	-- (0.1571,-0.5925)
	-- (0.1796,-0.5463)
	-- (0.2003,-0.4993)
	-- (0.2191,-0.4515)
	-- (0.2360,-0.4029)
	-- (0.2509,-0.3538)
	-- (0.2639,-0.3041)
	-- (0.2750,-0.2539)
	-- (0.2840,-0.2034)
	-- (0.2910,-0.1525)
	-- (0.2960,-0.1013)
	-- (0.2990,-0.0501)
	-- (0.3000,0.0013)
	-- (0.2989,0.0527)
	-- (0.2958,0.1039)
	-- (0.2907,0.1550)
	-- (0.2836,0.2059)
	-- (0.2744,0.2565)
	-- (0.2633,0.3066)
	-- (0.2502,0.3563)
	-- (0.2352,0.4054)
	-- (0.2182,0.4539)
	-- (0.1993,0.5017)
	-- (0.1785,0.5487)
	-- (0.1560,0.5948)
	-- (0.1316,0.6400)
	-- (0.1054,0.6842)
	-- (0.0775,0.7273)
	-- (0.0479,0.7693)
	-- (0.0167,0.8101)
	-- (-0.0161,0.8497)
	-- (-0.0504,0.8879)
	-- (-0.0863,0.9247)
	-- (-0.1235,0.9601)
	-- (-0.1621,0.9940)
	-- (-0.2020,1.0263)
	-- (-0.2432,1.0570)
	-- (-0.2856,1.0861)
	-- (-0.3290,1.1135)
	-- (-0.3735,1.1391)
	-- (-0.4190,1.1630)
	-- (-0.4654,1.1850)
	-- (-0.5127,1.2052)
	-- (-0.5607,1.2235)
	-- (-0.6093,1.2399)
	-- (-0.6586,1.2544)
	-- (-0.7085,1.2669)
	-- (-0.7587,1.2774)
	-- (-0.8094,1.2860)
	-- (-0.8603,1.2925)
	-- (-0.9115,1.2970)
	-- (-0.9628,1.2995)
	-- (-1.0142,1.2999)
	-- (-1.0655,1.2983)
	-- (-1.1168,1.2947)
	-- (-1.1678,1.2891)
	-- (-1.2186,1.2815)
	-- (-1.2691,1.2718)
}
\def\scaleA{0.6522}
\def\bildkurveB{
	(0.5400,-0.1500)
	-- (0.5609,-0.0417)
	-- (0.5723,0.0667)
	-- (0.5742,0.1747)
	-- (0.5667,0.2812)
	-- (0.5500,0.3857)
	-- (0.5244,0.4872)
	-- (0.4901,0.5852)
	-- (0.4476,0.6789)
	-- (0.3972,0.7676)
	-- (0.3394,0.8508)
	-- (0.2747,0.9278)
	-- (0.2037,0.9980)
	-- (0.1270,1.0611)
	-- (0.0453,1.1165)
	-- (-0.0409,1.1639)
	-- (-0.1307,1.2029)
	-- (-0.2235,1.2333)
	-- (-0.3186,1.2548)
	-- (-0.4150,1.2673)
	-- (-0.5122,1.2707)
	-- (-0.6093,1.2650)
	-- (-0.7055,1.2502)
	-- (-0.8000,1.2265)
	-- (-0.8921,1.1939)
	-- (-0.9811,1.1528)
	-- (-1.0662,1.1033)
	-- (-1.1468,1.0460)
	-- (-1.2221,0.9811)
	-- (-1.2916,0.9091)
	-- (-1.3546,0.8305)
	-- (-1.4106,0.7459)
	-- (-1.4590,0.6558)
	-- (-1.4995,0.5610)
	-- (-1.5316,0.4620)
	-- (-1.5550,0.3597)
	-- (-1.5694,0.2546)
	-- (-1.5746,0.1477)
	-- (-1.5703,0.0395)
	-- (-1.5566,-0.0690)
	-- (-1.5333,-0.1771)
	-- (-1.5004,-0.2839)
	-- (-1.4581,-0.3888)
	-- (-1.4066,-0.4908)
	-- (-1.3460,-0.5894)
	-- (-1.2766,-0.6836)
	-- (-1.1987,-0.7727)
	-- (-1.1128,-0.8562)
	-- (-1.0192,-0.9332)
	-- (-0.9186,-1.0032)
	-- (-0.8114,-1.0655)
	-- (-0.6983,-1.1197)
	-- (-0.5799,-1.1652)
	-- (-0.4569,-1.2015)
	-- (-0.3300,-1.2284)
	-- (-0.2000,-1.2454)
	-- (-0.0676,-1.2523)
	-- (0.0664,-1.2489)
	-- (0.2011,-1.2350)
	-- (0.3357,-1.2106)
	-- (0.4694,-1.1756)
	-- (0.6014,-1.1300)
	-- (0.7308,-1.0741)
	-- (0.8567,-1.0079)
	-- (0.9785,-0.9317)
	-- (1.0953,-0.8458)
	-- (1.2063,-0.7507)
	-- (1.3109,-0.6466)
	-- (1.4082,-0.5342)
	-- (1.4978,-0.4139)
	-- (1.5788,-0.2863)
	-- (1.6508,-0.1521)
	-- (1.7133,-0.0120)
	-- (1.7658,0.1333)
	-- (1.8078,0.2831)
	-- (1.8390,0.4366)
	-- (1.8592,0.5929)
	-- (1.8681,0.7513)
	-- (1.8655,0.9108)
	-- (1.8514,1.0706)
	-- (1.8256,1.2298)
	-- (1.7884,1.3876)
	-- (1.7397,1.5430)
	-- (1.6798,1.6953)
	-- (1.6088,1.8435)
	-- (1.5272,1.9869)
	-- (1.4352,2.1246)
	-- (1.3333,2.2559)
	-- (1.2219,2.3800)
	-- (1.1017,2.4963)
	-- (0.9732,2.6041)
	-- (0.8371,2.7027)
	-- (0.6940,2.7916)
	-- (0.5447,2.8704)
	-- (0.3899,2.9385)
	-- (0.2305,2.9956)
	-- (0.0673,3.0414)
	-- (-0.0989,3.0756)
	-- (-0.2672,3.0980)
	-- (-0.4367,3.1084)
	-- (-0.6066,3.1069)
	-- (-0.7759,3.0934)
	-- (-0.9437,3.0680)
	-- (-1.1093,3.0308)
	-- (-1.2716,2.9821)
	-- (-1.4299,2.9222)
	-- (-1.5833,2.8513)
	-- (-1.7311,2.7699)
	-- (-1.8725,2.6784)
	-- (-2.0068,2.5774)
	-- (-2.1332,2.4674)
	-- (-2.2512,2.3491)
	-- (-2.3602,2.2230)
	-- (-2.4596,2.0900)
	-- (-2.5490,1.9508)
	-- (-2.6280,1.8061)
	-- (-2.6961,1.6568)
	-- (-2.7532,1.5036)
	-- (-2.7990,1.3475)
	-- (-2.8333,1.1893)
	-- (-2.8561,1.0298)
	-- (-2.8672,0.8700)
	-- (-2.8669,0.7107)
	-- (-2.8551,0.5528)
	-- (-2.8321,0.3971)
	-- (-2.7981,0.2445)
	-- (-2.7533,0.0958)
	-- (-2.6983,-0.0483)
	-- (-2.6333,-0.1870)
	-- (-2.5590,-0.3196)
	-- (-2.4757,-0.4453)
	-- (-2.3841,-0.5637)
	-- (-2.2848,-0.6740)
	-- (-2.1786,-0.7758)
	-- (-2.0660,-0.8687)
	-- (-1.9479,-0.9521)
	-- (-1.8250,-1.0257)
	-- (-1.6980,-1.0893)
	-- (-1.5679,-1.1426)
	-- (-1.4354,-1.1855)
	-- (-1.3014,-1.2178)
	-- (-1.1667,-1.2396)
	-- (-1.0321,-1.2508)
	-- (-0.8985,-1.2516)
	-- (-0.7666,-1.2420)
	-- (-0.6373,-1.2225)
	-- (-0.5113,-1.1931)
	-- (-0.3894,-1.1544)
	-- (-0.2723,-1.1066)
	-- (-0.1606,-1.0503)
	-- (-0.0551,-0.9860)
	-- (0.0438,-0.9142)
	-- (0.1354,-0.8354)
	-- (0.2194,-0.7505)
	-- (0.2951,-0.6600)
	-- (0.3623,-0.5646)
	-- (0.4206,-0.4651)
	-- (0.4698,-0.3622)
	-- (0.5097,-0.2568)
	-- (0.5401,-0.1495)
	-- (0.5610,-0.0413)
	-- (0.5723,0.0672)
	-- (0.5741,0.1751)
	-- (0.5666,0.2817)
	-- (0.5499,0.3861)
	-- (0.5243,0.4876)
	-- (0.4900,0.5856)
	-- (0.4474,0.6793)
	-- (0.3970,0.7680)
	-- (0.3391,0.8511)
	-- (0.2744,0.9281)
	-- (0.2034,0.9983)
	-- (0.1267,1.0613)
	-- (0.0449,1.1167)
	-- (-0.0413,1.1641)
	-- (-0.1311,1.2030)
	-- (-0.2239,1.2334)
	-- (-0.3190,1.2548)
	-- (-0.4155,1.2673)
	-- (-0.5126,1.2707)
	-- (-0.6097,1.2649)
	-- (-0.7059,1.2501)
	-- (-0.8004,1.2263)
	-- (-0.8925,1.1937)
	-- (-0.9815,1.1526)
	-- (-1.0666,1.1031)
	-- (-1.1471,1.0457)
	-- (-1.2224,0.9808)
	-- (-1.2919,0.9087)
	-- (-1.3548,0.8301)
	-- (-1.4108,0.7455)
	-- (-1.4592,0.6554)
	-- (-1.4997,0.5606)
	-- (-1.5318,0.4616)
	-- (-1.5551,0.3592)
	-- (-1.5695,0.2542)
	-- (-1.5746,0.1472)
	-- (-1.5703,0.0391)
	-- (-1.5565,-0.0694)
	-- (-1.5331,-0.1775)
	-- (-1.5003,-0.2844)
	-- (-1.4579,-0.3892)
	-- (-1.4064,-0.4913)
	-- (-1.3457,-0.5898)
	-- (-1.2762,-0.6840)
	-- (-1.1983,-0.7731)
	-- (-1.1124,-0.8565)
	-- (-1.0188,-0.9335)
	-- (-0.9182,-1.0034)
	-- (-0.8110,-1.0657)
	-- (-0.6978,-1.1199)
	-- (-0.5794,-1.1653)
	-- (-0.4564,-1.2017)
	-- (-0.3295,-1.2285)
	-- (-0.1994,-1.2455)
	-- (-0.0670,-1.2524)
	-- (0.0670,-1.2489)
	-- (0.2017,-1.2350)
	-- (0.3363,-1.2105)
	-- (0.4700,-1.1754)
	-- (0.6019,-1.1298)
	-- (0.7313,-1.0738)
	-- (0.8573,-1.0076)
	-- (0.9790,-0.9314)
	-- (1.0958,-0.8455)
	-- (1.2068,-0.7503)
	-- (1.3113,-0.6462)
	-- (1.4086,-0.5337)
	-- (1.4981,-0.4134)
	-- (1.5791,-0.2858)
	-- (1.6511,-0.1516)
	-- (1.7135,-0.0114)
	-- (1.7660,0.1340)
	-- (1.8079,0.2838)
	-- (1.8391,0.4373)
	-- (1.8593,0.5936)
	-- (1.8681,0.7520)
	-- (1.8655,0.9115)
	-- (1.8513,1.0713)
	-- (1.8255,1.2305)
	-- (1.7882,1.3882)
	-- (1.7395,1.5437)
	-- (1.6795,1.6959)
	-- (1.6085,1.8441)
	-- (1.5268,1.9875)
	-- (1.4348,2.1252)
	-- (1.3328,2.2564)
	-- (1.2214,2.3805)
	-- (1.1012,2.4968)
	-- (0.9727,2.6045)
	-- (0.8365,2.7031)
	-- (0.6934,2.7920)
	-- (0.5440,2.8707)
	-- (0.3893,2.9388)
	-- (0.2298,2.9959)
	-- (0.0666,3.0416)
	-- (-0.0996,3.0757)
	-- (-0.2679,3.0980)
	-- (-0.4375,3.1084)
	-- (-0.6073,3.1068)
	-- (-0.7766,3.0933)
	-- (-0.9444,3.0678)
	-- (-1.1100,3.0306)
	-- (-1.2723,2.9819)
	-- (-1.4306,2.9219)
	-- (-1.5840,2.8510)
	-- (-1.7317,2.7695)
	-- (-1.8731,2.6780)
	-- (-2.0073,2.5770)
	-- (-2.1337,2.4669)
	-- (-2.2517,2.3486)
	-- (-2.3606,2.2225)
	-- (-2.4600,2.0895)
	-- (-2.5494,1.9502)
	-- (-2.6283,1.8055)
	-- (-2.6964,1.6561)
	-- (-2.7534,1.5030)
	-- (-2.7992,1.3468)
	-- (-2.8334,1.1886)
	-- (-2.8561,1.0292)
	-- (-2.8673,0.8693)
	-- (-2.8669,0.7101)
	-- (-2.8550,0.5522)
	-- (-2.8320,0.3965)
	-- (-2.7979,0.2439)
	-- (-2.7531,0.0952)
	-- (-2.6980,-0.0489)
	-- (-2.6330,-0.1876)
	-- (-2.5586,-0.3201)
	-- (-2.4753,-0.4458)
	-- (-2.3837,-0.5641)
	-- (-2.2844,-0.6745)
	-- (-2.1781,-0.7762)
	-- (-2.0655,-0.8690)
	-- (-1.9474,-0.9524)
	-- (-1.8244,-1.0260)
	-- (-1.6975,-1.0896)
	-- (-1.5674,-1.1429)
	-- (-1.4349,-1.1857)
	-- (-1.3008,-1.2180)
	-- (-1.1661,-1.2397)
	-- (-1.0315,-1.2508)
	-- (-0.8979,-1.2515)
	-- (-0.7660,-1.2420)
	-- (-0.6367,-1.2224)
	-- (-0.5108,-1.1930)
	-- (-0.3889,-1.1542)
	-- (-0.2718,-1.1064)
	-- (-0.1602,-1.0501)
	-- (-0.0546,-0.9857)
	-- (0.0442,-0.9138)
	-- (0.1358,-0.8351)
	-- (0.2197,-0.7501)
	-- (0.2954,-0.6596)
	-- (0.3626,-0.5642)
	-- (0.4209,-0.4647)
	-- (0.4700,-0.3618)
	-- (0.5098,-0.2563)
	-- (0.5402,-0.1491)
	-- (0.5610,-0.0408)
	-- (0.5723,0.0677)
	-- (0.5741,0.1756)
	-- (0.5666,0.2821)
	-- (0.5498,0.3865)
	-- (0.5241,0.4881)
	-- (0.4898,0.5860)
	-- (0.4472,0.6797)
	-- (0.3967,0.7684)
	-- (0.3389,0.8515)
	-- (0.2742,0.9284)
	-- (0.2031,0.9986)
	-- (0.1264,1.0616)
	-- (0.0446,1.1169)
	-- (-0.0416,1.1643)
	-- (-0.1315,1.2032)
	-- (-0.2243,1.2335)
	-- (-0.3194,1.2549)
	-- (-0.4159,1.2673)
	-- (-0.5130,1.2707)
	-- (-0.6101,1.2649)
	-- (-0.7063,1.2500)
	-- (-0.8008,1.2262)
	-- (-0.8929,1.1936)
	-- (-0.9819,1.1524)
	-- (-1.0669,1.1029)
	-- (-1.1475,1.0455)
	-- (-1.2227,0.9805)
	-- (-1.2921,0.9084)
	-- (-1.3551,0.8298)
	-- (-1.4110,0.7451)
	-- (-1.4594,0.6550)
	-- (-1.4998,0.5602)
	-- (-1.5319,0.4612)
	-- (-1.5552,0.3588)
	-- (-1.5695,0.2537)
	-- (-1.5746,0.1467)
	-- (-1.5702,0.0386)
	-- (-1.5564,-0.0699)
	-- (-1.5330,-0.1780)
	-- (-1.5001,-0.2848)
	-- (-1.4578,-0.3896)
	-- (-1.4061,-0.4917)
	-- (-1.3454,-0.5902)
	-- (-1.2759,-0.6843)
	-- (-1.1980,-0.7735)
	-- (-1.1120,-0.8568)
	-- (-1.0184,-0.9338)
	-- (-0.9177,-1.0037)
	-- (-0.8105,-1.0660)
	-- (-0.6974,-1.1201)
	-- (-0.5789,-1.1655)
	-- (-0.4559,-1.2018)
	-- (-0.3289,-1.2286)
	-- (-0.1989,-1.2455)
	-- (-0.0664,-1.2524)
	-- (0.0675,-1.2489)
	-- (0.2022,-1.2349)
	-- (0.3369,-1.2103)
	-- (0.4705,-1.1752)
	-- (0.6025,-1.1296)
	-- (0.7318,-1.0736)
	-- (0.8578,-1.0073)
	-- (0.9795,-0.9310)
	-- (1.0963,-0.8451)
	-- (1.2072,-0.7498)
	-- (1.3117,-0.6457)
	-- (1.4090,-0.5332)
	-- (1.4985,-0.4129)
	-- (1.5795,-0.2852)
	-- (1.6514,-0.1510)
	-- (1.7138,-0.0108)
	-- (1.7662,0.1346)
	-- (1.8081,0.2844)
	-- (1.8392,0.4379)
	-- (1.8593,0.5943)
	-- (1.8681,0.7526)
	-- (1.8654,0.9121)
	-- (1.8512,1.0719)
	-- (1.8254,1.2312)
	-- (1.7880,1.3889)
	-- (1.7393,1.5443)
	-- (1.6792,1.6966)
	-- (1.6082,1.8447)
	-- (1.5265,1.9881)
	-- (1.4344,2.1257)
	-- (1.3324,2.2570)
	-- (1.2210,2.3810)
	-- (1.1007,2.4973)
	-- (0.9721,2.6049)
	-- (0.8359,2.7035)
	-- (0.6928,2.7923)
	-- (0.5434,2.8710)
	-- (0.3886,2.9391)
	-- (0.2291,2.9961)
	-- (0.0659,3.0417)
	-- (-0.1003,3.0758)
	-- (-0.2686,3.0981)
	-- (-0.4382,3.1084)
	-- (-0.6080,3.1068)
	-- (-0.7773,3.0932)
	-- (-0.9451,3.0677)
	-- (-1.1106,3.0305)
	-- (-1.2729,2.9817)
	-- (-1.4312,2.9216)
	-- (-1.5846,2.8506)
	-- (-1.7324,2.7691)
	-- (-1.8737,2.6776)
	-- (-2.0079,2.5765)
	-- (-2.1343,2.4664)
	-- (-2.2522,2.3480)
	-- (-2.3611,2.2220)
	-- (-2.4604,2.0889)
	-- (-2.5497,1.9496)
	-- (-2.6286,1.8049)
	-- (-2.6967,1.6555)
	-- (-2.7537,1.5023)
	-- (-2.7993,1.3462)
	-- (-2.8335,1.1879)
	-- (-2.8562,1.0285)
	-- (-2.8673,0.8687)
	-- (-2.8668,0.7094)
	-- (-2.8550,0.5515)
	-- (-2.8318,0.3958)
	-- (-2.7977,0.2432)
	-- (-2.7529,0.0945)
	-- (-2.6978,-0.0495)
	-- (-2.6327,-0.1881)
	-- (-2.5583,-0.3206)
	-- (-2.4749,-0.4463)
	-- (-2.3833,-0.5646)
	-- (-2.2840,-0.6749)
	-- (-2.1777,-0.7767)
	-- (-2.0650,-0.8694)
	-- (-1.9469,-0.9527)
	-- (-1.8239,-1.0263)
	-- (-1.6969,-1.0898)
	-- (-1.5668,-1.1431)
	-- (-1.4343,-1.1858)
	-- (-1.3003,-1.2181)
	-- (-1.1656,-1.2397)
	-- (-1.0310,-1.2508)
	-- (-0.8973,-1.2515)
	-- (-0.7655,-1.2419)
	-- (-0.6362,-1.2223)
	-- (-0.5102,-1.1929)
	-- (-0.3884,-1.1540)
	-- (-0.2713,-1.1062)
	-- (-0.1597,-1.0498)
	-- (-0.0542,-0.9854)
	-- (0.0446,-0.9135)
	-- (0.1362,-0.8348)
	-- (0.2200,-0.7497)
	-- (0.2957,-0.6592)
	-- (0.3628,-0.5638)
	-- (0.4211,-0.4642)
	-- (0.4702,-0.3614)
	-- (0.5100,-0.2559)
	-- (0.5403,-0.1486)
	-- (0.5611,-0.0404)
	-- (0.5724,0.0681)
	-- (0.5741,0.1760)
	-- (0.5665,0.2825)
	-- (0.5497,0.3870)
	-- (0.5240,0.4885)
	-- (0.4896,0.5864)
	-- (0.4470,0.6801)
	-- (0.3965,0.7687)
	-- (0.3386,0.8518)
	-- (0.2739,0.9287)
	-- (0.2028,0.9989)
	-- (0.1260,1.0618)
	-- (0.0442,1.1172)
	-- (-0.0420,1.1644)
	-- (-0.1319,1.2033)
	-- (-0.2247,1.2336)
	-- (-0.3198,1.2550)
	-- (-0.4163,1.2674)
	-- (-0.5134,1.2707)
	-- (-0.6105,1.2649)
	-- (-0.7067,1.2500)
	-- (-0.8012,1.2261)
	-- (-0.8933,1.1934)
	-- (-0.9822,1.1522)
	-- (-1.0673,1.1027)
	-- (-1.1478,1.0452)
	-- (-1.2230,0.9802)
	-- (-1.2924,0.9081)
	-- (-1.3553,0.8294)
	-- (-1.4112,0.7448)
	-- (-1.4596,0.6547)
	-- (-1.5000,0.5598)
	-- (-1.5320,0.4608)
	-- (-1.5553,0.3584)
	-- (-1.5695,0.2533)
	-- (-1.5746,0.1463)
	-- (-1.5702,0.0382)
	-- (-1.5563,-0.0703)
	-- (-1.5329,-0.1784)
	-- (-1.4999,-0.2853)
	-- (-1.4576,-0.3901)
	-- (-1.4059,-0.4921)
	-- (-1.3452,-0.5906)
	-- (-1.2756,-0.6847)
	-- (-1.1977,-0.7738)
	-- (-1.1116,-0.8572)
	-- (-1.0180,-0.9341)
	-- (-0.9173,-1.0040)
	-- (-0.8100,-1.0662)
	-- (-0.6969,-1.1203)
	-- (-0.5784,-1.1657)
	-- (-0.4553,-1.2019)
	-- (-0.3284,-1.2287)
	-- (-0.1983,-1.2456)
	-- (-0.0659,-1.2524)
	-- (0.0681,-1.2488)
	-- (0.2028,-1.2348)
	-- (0.3374,-1.2102)
	-- (0.4711,-1.1751)
	-- (0.6030,-1.1294)
	-- (0.7324,-1.0733)
	-- (0.8583,-1.0070)
	-- (0.9800,-0.9307)
	-- (1.0967,-0.8447)
	-- (1.2077,-0.7494)
	-- (1.3122,-0.6453)
	-- (1.4094,-0.5327)
	-- (1.4988,-0.4123)
	-- (1.5798,-0.2847)
	-- (1.6517,-0.1504)
	-- (1.7140,-0.0102)
	-- (1.7664,0.1352)
	-- (1.8082,0.2851)
	-- (1.8393,0.4386)
	-- (1.8594,0.5949)
	-- (1.8681,0.7533)
	-- (1.8654,0.9128)
	-- (1.8511,1.0726)
	-- (1.8252,1.2318)
	-- (1.7879,1.3896)
	-- (1.7390,1.5450)
	-- (1.6790,1.6972)
	-- (1.6079,1.8454)
	-- (1.5261,1.9887)
	-- (1.4340,2.1263)
	-- (1.3319,2.2575)
	-- (1.2205,2.3816)
	-- (1.1001,2.4977)
	-- (0.9715,2.6054)
	-- (0.8353,2.7039)
	-- (0.6921,2.7927)
	-- (0.5428,2.8713)
	-- (0.3879,2.9393)
	-- (0.2285,2.9963)
	-- (0.0652,3.0419)
	-- (-0.1010,3.0759)
	-- (-0.2694,3.0982)
	-- (-0.4389,3.1085)
	-- (-0.6087,3.1068)
	-- (-0.7780,3.0931)
	-- (-0.9458,3.0676)
	-- (-1.1113,3.0303)
	-- (-1.2736,2.9814)
	-- (-1.4319,2.9213)
	-- (-1.5852,2.8503)
	-- (-1.7330,2.7688)
	-- (-1.8743,2.6772)
	-- (-2.0084,2.5761)
	-- (-2.1348,2.4660)
	-- (-2.2527,2.3475)
	-- (-2.3615,2.2214)
	-- (-2.4608,2.0883)
	-- (-2.5501,1.9490)
	-- (-2.6289,1.8042)
	-- (-2.6969,1.6549)
	-- (-2.7539,1.5017)
	-- (-2.7995,1.3455)
	-- (-2.8337,1.1873)
	-- (-2.8563,1.0278)
	-- (-2.8673,0.8680)
	-- (-2.8668,0.7087)
	-- (-2.8549,0.5508)
	-- (-2.8317,0.3952)
	-- (-2.7976,0.2426)
	-- (-2.7527,0.0939)
	-- (-2.6975,-0.0501)
	-- (-2.6325,-0.1887)
	-- (-2.5580,-0.3212)
	-- (-2.4746,-0.4469)
	-- (-2.3829,-0.5651)
	-- (-2.2835,-0.6754)
	-- (-2.1772,-0.7771)
	-- (-2.0645,-0.8698)
	-- (-1.9464,-0.9531)
	-- (-1.8234,-1.0266)
	-- (-1.6964,-1.0901)
	-- (-1.5663,-1.1433)
	-- (-1.4337,-1.1860)
	-- (-1.2997,-1.2182)
	-- (-1.1650,-1.2398)
	-- (-1.0304,-1.2509)
	-- (-0.8968,-1.2515)
	-- (-0.7649,-1.2419)
	-- (-0.6357,-1.2222)
	-- (-0.5097,-1.1927)
	-- (-0.3879,-1.1539)
	-- (-0.2708,-1.1060)
	-- (-0.1592,-1.0496)
	-- (-0.0538,-0.9851)
	-- (0.0450,-0.9132)
	-- (0.1366,-0.8344)
	-- (0.2204,-0.7494)
	-- (0.2960,-0.6588)
	-- (0.3631,-0.5634)
	-- (0.4213,-0.4638)
	-- (0.4704,-0.3609)
	-- (0.5101,-0.2555)
	-- (0.5404,-0.1482)
	-- (0.5612,-0.0399)
	-- (0.5724,0.0686)
	-- (0.5741,0.1765)
	-- (0.5665,0.2830)
	-- (0.5496,0.3874)
	-- (0.5239,0.4889)
	-- (0.4895,0.5868)
	-- (0.4468,0.6804)
	-- (0.3963,0.7691)
	-- (0.3384,0.8521)
	-- (0.2736,0.9290)
	-- (0.2025,0.9992)
	-- (0.1257,1.0621)
	-- (0.0439,1.1174)
	-- (-0.0424,1.1646)
	-- (-0.1323,1.2035)
	-- (-0.2251,1.2337)
	-- (-0.3202,1.2551)
	-- (-0.4167,1.2674)
	-- (-0.5138,1.2707)
	-- (-0.6109,1.2648)
	-- (-0.7071,1.2499)
	-- (-0.8016,1.2260)
	-- (-0.8937,1.1933)
	-- (-0.9826,1.1520)
	-- (-1.0676,1.1024)
	-- (-1.1481,1.0449)
	-- (-1.2233,0.9799)
	-- (-1.2927,0.9078)
	-- (-1.3556,0.8291)
	-- (-1.4114,0.7444)
	-- (-1.4598,0.6543)
	-- (-1.5001,0.5594)
	-- (-1.5321,0.4603)
	-- (-1.5553,0.3579)
	-- (-1.5696,0.2528)
	-- (-1.5746,0.1458)
	-- (-1.5702,0.0377)
	-- (-1.5562,-0.0708)
	-- (-1.5328,-0.1789)
	-- (-1.4998,-0.2857)
	-- (-1.4574,-0.3905)
	-- (-1.4056,-0.4925)
	-- (-1.3449,-0.5910)
	-- (-1.2753,-0.6851)
	-- (-1.1973,-0.7742)
	-- (-1.1113,-0.8575)
	-- (-1.0176,-0.9344)
	-- (-0.9169,-1.0043)
	-- (-0.8096,-1.0665)
	-- (-0.6964,-1.1205)
	-- (-0.5779,-1.1659)
	-- (-0.4548,-1.2021)
	-- (-0.3278,-1.2288)
	-- (-0.1977,-1.2456)
	-- (-0.0653,-1.2524)
	-- (0.0687,-1.2488)
	-- (0.2034,-1.2347)
	-- (0.3380,-1.2101)
	-- (0.4717,-1.1749)
	-- (0.6036,-1.1292)
	-- (0.7329,-1.0730)
	-- (0.8588,-1.0067)
	-- (0.9805,-0.9303)
	-- (1.0972,-0.8443)
	-- (1.2081,-0.7490)
	-- (1.3126,-0.6448)
	-- (1.4098,-0.5322)
	-- (1.4992,-0.4118)
	-- (1.5801,-0.2841)
	-- (1.6520,-0.1498)
	-- (1.7143,-0.0096)
	-- (1.7666,0.1358)
	-- (1.8084,0.2857)
	-- (1.8395,0.4392)
	-- (1.8594,0.5956)
	-- (1.8681,0.7540)
	-- (1.8653,0.9135)
	-- (1.8510,1.0733)
	-- (1.8251,1.2325)
	-- (1.7877,1.3902)
	-- (1.7388,1.5456)
	-- (1.6787,1.6978)
	-- (1.6076,1.8460)
	-- (1.5257,1.9893)
	-- (1.4335,2.1269)
	-- (1.3315,2.2581)
	-- (1.2200,2.3821)
	-- (1.0996,2.4982)
	-- (0.9710,2.6058)
	-- (0.8347,2.7043)
	-- (0.6915,2.7930)
	-- (0.5421,2.8716)
	-- (0.3873,2.9396)
	-- (0.2278,2.9965)
	-- (0.0645,3.0421)
	-- (-0.1017,3.0761)
	-- (-0.2701,3.0982)
	-- (-0.4396,3.1085)
	-- (-0.6095,3.1067)
	-- (-0.7787,3.0930)
	-- (-0.9465,3.0674)
	-- (-1.1120,3.0301)
	-- (-1.2743,2.9812)
	-- (-1.4325,2.9211)
	-- (-1.5859,2.8500)
	-- (-1.7336,2.7684)
	-- (-1.8748,2.6768)
	-- (-2.0090,2.5756)
	-- (-2.1353,2.4655)
	-- (-2.2531,2.3470)
	-- (-2.3620,2.2209)
	-- (-2.4612,2.0877)
	-- (-2.5504,1.9484)
	-- (-2.6292,1.8036)
	-- (-2.6972,1.6542)
	-- (-2.7541,1.5010)
	-- (-2.7997,1.3448)
	-- (-2.8338,1.1866)
	-- (-2.8563,1.0271)
	-- (-2.8673,0.8673)
	-- (-2.8668,0.7081)
	-- (-2.8548,0.5502)
	-- (-2.8316,0.3945)
	-- (-2.7974,0.2420)
	-- (-2.7525,0.0933)
	-- (-2.6973,-0.0507)
	-- (-2.6322,-0.1893)
	-- (-2.5576,-0.3217)
	-- (-2.4742,-0.4474)
	-- (-2.3825,-0.5656)
	-- (-2.2831,-0.6758)
	-- (-2.1767,-0.7775)
	-- (-2.0641,-0.8702)
	-- (-1.9458,-0.9534)
	-- (-1.8228,-1.0269)
	-- (-1.6959,-1.0903)
	-- (-1.5657,-1.1435)
	-- (-1.4332,-1.1861)
	-- (-1.2991,-1.2183)
	-- (-1.1644,-1.2399)
	-- (-1.0298,-1.2509)
	-- (-0.8962,-1.2515)
	-- (-0.7644,-1.2418)
	-- (-0.6351,-1.2221)
	-- (-0.5092,-1.1926)
	-- (-0.3874,-1.1537)
	-- (-0.2703,-1.1058)
	-- (-0.1588,-1.0493)
	-- (-0.0533,-0.9848)
	-- (0.0454,-0.9129)
	-- (0.1369,-0.8341)
	-- (0.2207,-0.7490)
	-- (0.2963,-0.6584)
	-- (0.3634,-0.5629)
	-- (0.4215,-0.4634)
	-- (0.4706,-0.3605)
	-- (0.5103,-0.2550)
	-- (0.5405,-0.1477)
	-- (0.5612,-0.0395)
	-- (0.5724,0.0690)
	-- (0.5741,0.1769)
	-- (0.5664,0.2834)
	-- (0.5495,0.3878)
	-- (0.5237,0.4893)
	-- (0.4893,0.5872)
	-- (0.4466,0.6808)
	-- (0.3960,0.7694)
	-- (0.3381,0.8525)
	-- (0.2733,0.9293)
	-- (0.2022,0.9994)
	-- (0.1254,1.0623)
	-- (0.0435,1.1176)
	-- (-0.0427,1.1648)
	-- (-0.1327,1.2036)
	-- (-0.2255,1.2338)
	-- (-0.3206,1.2551)
	-- (-0.4171,1.2674)
	-- (-0.5143,1.2707)
	-- (-0.6113,1.2648)
	-- (-0.7075,1.2498)
	-- (-0.8020,1.2259)
	-- (-0.8940,1.1931)
	-- (-0.9830,1.1518)
	-- (-1.0680,1.1022)
	-- (-1.1484,1.0447)
	-- (-1.2237,0.9796)
	-- (-1.2930,0.9075)
	-- (-1.3558,0.8287)
	-- (-1.4117,0.7440)
	-- (-1.4600,0.6539)
	-- (-1.5003,0.5589)
	-- (-1.5322,0.4599)
	-- (-1.5554,0.3575)
	-- (-1.5696,0.2524)
	-- (-1.5746,0.1454)
	-- (-1.5701,0.0372)
	-- (-1.5562,-0.0713)
	-- (-1.5327,-0.1793)
	-- (-1.4996,-0.2861)
	-- (-1.4572,-0.3910)
	-- (-1.4054,-0.4930)
	-- (-1.3446,-0.5914)
	-- (-1.2750,-0.6855)
	-- (-1.1970,-0.7745)
	-- (-1.1109,-0.8578)
	-- (-1.0172,-0.9347)
	-- (-0.9164,-1.0045)
	-- (-0.8091,-1.0667)
	-- (-0.6959,-1.1207)
	-- (-0.5774,-1.1660)
	-- (-0.4543,-1.2022)
	-- (-0.3273,-1.2289)
	-- (-0.1972,-1.2457)
	-- (-0.0648,-1.2524)
	-- (0.0692,-1.2487)
	-- (0.2039,-1.2346)
	-- (0.3386,-1.2100)
	-- (0.4722,-1.1747)
	-- (0.6041,-1.1290)
	-- (0.7335,-1.0728)
	-- (0.8594,-1.0064)
	-- (0.9810,-0.9300)
	-- (1.0977,-0.8439)
	-- (1.2086,-0.7486)
	-- (1.3130,-0.6444)
	-- (1.4102,-0.5317)
	-- (1.4995,-0.4113)
	-- (1.5804,-0.2836)
	-- (1.6522,-0.1493)
	-- (1.7145,-0.0090)
	-- (1.7667,0.1364)
	-- (1.8085,0.2863)
	-- (1.8396,0.4399)
	-- (1.8595,0.5963)
	-- (1.8681,0.7546)
	-- (1.8653,0.9142)
	-- (1.8509,1.0740)
	-- (1.8250,1.2332)
	-- (1.7875,1.3909)
	-- (1.7386,1.5463)
	-- (1.6784,1.6984)
	-- (1.6072,1.8466)
	-- (1.5254,1.9899)
	-- (1.4331,2.1274)
	-- (1.3310,2.2586)
	-- (1.2195,2.3826)
	-- (1.0991,2.4987)
	-- (0.9704,2.6062)
	-- (0.8341,2.7047)
	-- (0.6909,2.7934)
	-- (0.5415,2.8719)
	-- (0.3866,2.9398)
	-- (0.2271,2.9967)
	-- (0.0638,3.0423)
	-- (-0.1024,3.0762)
	-- (-0.2708,3.0983)
	-- (-0.4403,3.1085)
	-- (-0.6102,3.1067)
	-- (-0.7794,3.0930)
	-- (-0.9473,3.0673)
	-- (-1.1127,3.0299)
	-- (-1.2750,2.9810)
	-- (-1.4332,2.9208)
	-- (-1.5865,2.8497)
	-- (-1.7342,2.7681)
	-- (-1.8754,2.6764)
	-- (-2.0095,2.5752)
	-- (-2.1358,2.4650)
	-- (-2.2536,2.3465)
	-- (-2.3624,2.2203)
	-- (-2.4616,2.0872)
	-- (-2.5508,1.9478)
	-- (-2.6295,1.8030)
	-- (-2.6975,1.6536)
	-- (-2.7543,1.5004)
	-- (-2.7998,1.3442)
	-- (-2.8339,1.1859)
	-- (-2.8564,1.0265)
	-- (-2.8674,0.8667)
	-- (-2.8668,0.7074)
	-- (-2.8547,0.5495)
	-- (-2.8315,0.3939)
	-- (-2.7972,0.2413)
	-- (-2.7523,0.0927)
	-- (-2.6970,-0.0513)
	-- (-2.6319,-0.1898)
	-- (-2.5573,-0.3223)
	-- (-2.4738,-0.4479)
	-- (-2.3821,-0.5661)
	-- (-2.2827,-0.6762)
	-- (-2.1763,-0.7779)
	-- (-2.0636,-0.8705)
	-- (-1.9453,-0.9537)
	-- (-1.8223,-1.0272)
	-- (-1.6953,-1.0906)
	-- (-1.5651,-1.1437)
	-- (-1.4326,-1.1863)
	-- (-1.2986,-1.2184)
	-- (-1.1639,-1.2399)
	-- (-1.0293,-1.2509)
	-- (-0.8957,-1.2515)
	-- (-0.7638,-1.2417)
	-- (-0.6346,-1.2220)
	-- (-0.5087,-1.1924)
	-- (-0.3869,-1.1535)
	-- (-0.2699,-1.1055)
	-- (-0.1583,-1.0491)
	-- (-0.0529,-0.9846)
	-- (0.0458,-0.9126)
	-- (0.1373,-0.8337)
	-- (0.2210,-0.7486)
	-- (0.2966,-0.6580)
	-- (0.3636,-0.5625)
	-- (0.4218,-0.4630)
	-- (0.4708,-0.3600)
	-- (0.5104,-0.2546)
	-- (0.5406,-0.1473)
	-- (0.5613,-0.0390)
	-- (0.5724,0.0695)
	-- (0.5741,0.1774)
	-- (0.5664,0.2839)
	-- (0.5495,0.3883)
	-- (0.5236,0.4898)
	-- (0.4891,0.5876)
	-- (0.4464,0.6812)
	-- (0.3958,0.7698)
	-- (0.3378,0.8528)
	-- (0.2730,0.9296)
	-- (0.2019,0.9997)
	-- (0.1250,1.0626)
	-- (0.0432,1.1178)
	-- (-0.0431,1.1650)
	-- (-0.1330,1.2038)
	-- (-0.2259,1.2339)
	-- (-0.3210,1.2552)
	-- (-0.4175,1.2675)
	-- (-0.5147,1.2707)
	-- (-0.6117,1.2647)
	-- (-0.7079,1.2497)
	-- (-0.8024,1.2257)
	-- (-0.8944,1.1930)
	-- (-0.9833,1.1516)
	-- (-1.0683,1.1020)
	-- (-1.1488,1.0444)
	-- (-1.2240,0.9793)
	-- (-1.2933,0.9071)
	-- (-1.3561,0.8284)
	-- (-1.4119,0.7436)
	-- (-1.4601,0.6535)
	-- (-1.5004,0.5585)
	-- (-1.5323,0.4595)
	-- (-1.5555,0.3570)
	-- (-1.5697,0.2519)
	-- (-1.5746,0.1449)
	-- (-1.5701,0.0368)
	-- (-1.5561,-0.0717)
	-- (-1.5325,-0.1798)
	-- (-1.4995,-0.2866)
	-- (-1.4570,-0.3914)
	-- (-1.4052,-0.4934)
	-- (-1.3443,-0.5918)
	-- (-1.2747,-0.6859)
}
\def\scaleB{0.4590}
\def\bildkurveC{
	(0.9150,0.1530)
	-- (0.8722,0.3589)
	-- (0.8025,0.5592)
	-- (0.7069,0.7506)
	-- (0.5864,0.9299)
	-- (0.4427,1.0941)
	-- (0.2777,1.2403)
	-- (0.0935,1.3660)
	-- (-0.1071,1.4688)
	-- (-0.3214,1.5468)
	-- (-0.5463,1.5984)
	-- (-0.7785,1.6222)
	-- (-1.0147,1.6175)
	-- (-1.2514,1.5837)
	-- (-1.4851,1.5209)
	-- (-1.7123,1.4294)
	-- (-1.9294,1.3101)
	-- (-2.1333,1.1642)
	-- (-2.3207,0.9932)
	-- (-2.4886,0.7993)
	-- (-2.6343,0.5847)
	-- (-2.7553,0.3521)
	-- (-2.8496,0.1044)
	-- (-2.9154,-0.1552)
	-- (-2.9513,-0.4232)
	-- (-2.9564,-0.6963)
	-- (-2.9300,-0.9707)
	-- (-2.8721,-1.2427)
	-- (-2.7830,-1.5087)
	-- (-2.6634,-1.7651)
	-- (-2.5145,-2.0083)
	-- (-2.3379,-2.2348)
	-- (-2.1356,-2.4416)
	-- (-1.9099,-2.6257)
	-- (-1.6634,-2.7844)
	-- (-1.3993,-2.9154)
	-- (-1.1207,-3.0167)
	-- (-0.8310,-3.0867)
	-- (-0.5340,-3.1243)
	-- (-0.2333,-3.1287)
	-- (0.0672,-3.0995)
	-- (0.3637,-3.0370)
	-- (0.6523,-2.9417)
	-- (0.9294,-2.8145)
	-- (1.1913,-2.6571)
	-- (1.4347,-2.4712)
	-- (1.6563,-2.2590)
	-- (1.8532,-2.0232)
	-- (2.0230,-1.7667)
	-- (2.1633,-1.4927)
	-- (2.2723,-1.2047)
	-- (2.3486,-0.9063)
	-- (2.3912,-0.6012)
	-- (2.3994,-0.2935)
	-- (2.3732,0.0131)
	-- (2.3130,0.3146)
	-- (2.2194,0.6070)
	-- (2.0939,0.8868)
	-- (1.9379,1.1503)
	-- (1.7535,1.3941)
	-- (1.5433,1.6151)
	-- (1.3100,1.8106)
	-- (1.0566,1.9781)
	-- (0.7865,2.1154)
	-- (0.5033,2.2210)
	-- (0.2107,2.2934)
	-- (-0.0874,2.3320)
	-- (-0.3872,2.3363)
	-- (-0.6846,2.3064)
	-- (-0.9758,2.2429)
	-- (-1.2569,2.1466)
	-- (-1.5244,2.0191)
	-- (-1.7746,1.8622)
	-- (-2.0044,1.6780)
	-- (-2.2108,1.4692)
	-- (-2.3912,1.2387)
	-- (-2.5433,0.9896)
	-- (-2.6652,0.7253)
	-- (-2.7554,0.4495)
	-- (-2.8128,0.1659)
	-- (-2.8369,-0.1216)
	-- (-2.8274,-0.4092)
	-- (-2.7847,-0.6930)
	-- (-2.7096,-0.9691)
	-- (-2.6030,-1.2340)
	-- (-2.4668,-1.4840)
	-- (-2.3028,-1.7159)
	-- (-2.1134,-1.9266)
	-- (-1.9012,-2.1134)
	-- (-1.6694,-2.2739)
	-- (-1.4210,-2.4060)
	-- (-1.1596,-2.5080)
	-- (-0.8887,-2.5788)
	-- (-0.6121,-2.6176)
	-- (-0.3335,-2.6239)
	-- (-0.0568,-2.5980)
	-- (0.2145,-2.5402)
	-- (0.4765,-2.4516)
	-- (0.7258,-2.3334)
	-- (0.9591,-2.1876)
	-- (1.1734,-2.0162)
	-- (1.3658,-1.8217)
	-- (1.5339,-1.6069)
	-- (1.6756,-1.3749)
	-- (1.7891,-1.1288)
	-- (1.8731,-0.8721)
	-- (1.9267,-0.6084)
	-- (1.9493,-0.3413)
	-- (1.9409,-0.0744)
	-- (1.9019,0.1886)
	-- (1.8330,0.4441)
	-- (1.7354,0.6888)
	-- (1.6108,0.9194)
	-- (1.4611,1.1329)
	-- (1.2885,1.3265)
	-- (1.0957,1.4978)
	-- (0.8856,1.6446)
	-- (0.6612,1.7652)
	-- (0.4258,1.8582)
	-- (0.1828,1.9226)
	-- (-0.0644,1.9578)
	-- (-0.3122,1.9637)
	-- (-0.5571,1.9406)
	-- (-0.7957,1.8891)
	-- (-1.0247,1.8103)
	-- (-1.2410,1.7058)
	-- (-1.4415,1.5773)
	-- (-1.6237,1.4270)
	-- (-1.7851,1.2574)
	-- (-1.9237,1.0712)
	-- (-2.0377,0.8714)
	-- (-2.1257,0.6612)
	-- (-2.1868,0.4437)
	-- (-2.2205,0.2223)
	-- (-2.2265,0.0005)
	-- (-2.2052,-0.2184)
	-- (-2.1571,-0.4311)
	-- (-2.0835,-0.6344)
	-- (-1.9856,-0.8252)
	-- (-1.8654,-1.0006)
	-- (-1.7249,-1.1582)
	-- (-1.5665,-1.2955)
	-- (-1.3930,-1.4106)
	-- (-1.2073,-1.5018)
	-- (-1.0123,-1.5678)
	-- (-0.8115,-1.6077)
	-- (-0.6080,-1.6210)
	-- (-0.4052,-1.6076)
	-- (-0.2065,-1.5679)
	-- (-0.0150,-1.5024)
	-- (0.1660,-1.4125)
	-- (0.3335,-1.2994)
	-- (0.4848,-1.1652)
	-- (0.6174,-1.0119)
	-- (0.7289,-0.8420)
	-- (0.8174,-0.6582)
	-- (0.8813,-0.4636)
	-- (0.9193,-0.2611)
	-- (0.9307,-0.0542)
	-- (0.9149,0.1539)
	-- (0.8719,0.3597)
	-- (0.8022,0.5600)
	-- (0.7064,0.7514)
	-- (0.5859,0.9306)
	-- (0.4421,1.0947)
	-- (0.2769,1.2409)
	-- (0.0927,1.3665)
	-- (-0.1080,1.4692)
	-- (-0.3223,1.5471)
	-- (-0.5472,1.5985)
	-- (-0.7795,1.6222)
	-- (-1.0157,1.6174)
	-- (-1.2524,1.5835)
	-- (-1.4861,1.5205)
	-- (-1.7132,1.4290)
	-- (-1.9303,1.3095)
	-- (-2.1341,1.1635)
	-- (-2.3214,0.9924)
	-- (-2.4892,0.7984)
	-- (-2.6348,0.5837)
	-- (-2.7558,0.3511)
	-- (-2.8500,0.1033)
	-- (-2.9156,-0.1563)
	-- (-2.9514,-0.4244)
	-- (-2.9563,-0.6974)
	-- (-2.9298,-0.9718)
	-- (-2.8718,-1.2439)
	-- (-2.7825,-1.5099)
	-- (-2.6628,-1.7662)
	-- (-2.5138,-2.0093)
	-- (-2.3371,-2.2357)
	-- (-2.1347,-2.4425)
	-- (-1.9089,-2.6264)
	-- (-1.6624,-2.7850)
	-- (-1.3982,-2.9159)
	-- (-1.1195,-3.0171)
	-- (-0.8298,-3.0870)
	-- (-0.5327,-3.1244)
	-- (-0.2320,-3.1286)
	-- (0.0685,-3.0993)
	-- (0.3649,-3.0367)
	-- (0.6535,-2.9412)
	-- (0.9305,-2.8139)
	-- (1.1924,-2.6564)
	-- (1.4356,-2.4703)
	-- (1.6571,-2.2580)
	-- (1.8540,-2.0222)
	-- (2.0237,-1.7656)
	-- (2.1638,-1.4916)
	-- (2.2727,-1.2035)
	-- (2.3489,-0.9050)
	-- (2.3913,-0.5999)
	-- (2.3994,-0.2922)
	-- (2.3731,0.0144)
	-- (2.3127,0.3158)
	-- (2.2190,0.6083)
	-- (2.0933,0.8880)
	-- (1.9371,1.1513)
	-- (1.7527,1.3951)
	-- (1.5424,1.6160)
	-- (1.3090,1.8114)
	-- (1.0555,1.9787)
	-- (0.7854,2.1159)
	-- (0.5021,2.2213)
	-- (0.2095,2.2937)
	-- (-0.0887,2.3321)
	-- (-0.3884,2.3363)
	-- (-0.6858,2.3062)
	-- (-0.9770,2.2425)
	-- (-1.2581,2.1461)
	-- (-1.5254,2.0185)
	-- (-1.7756,1.8615)
	-- (-2.0053,1.6772)
	-- (-2.2117,1.4683)
	-- (-2.3919,1.2377)
	-- (-2.5439,0.9885)
	-- (-2.6657,0.7242)
	-- (-2.7557,0.4483)
	-- (-2.8130,0.1647)
	-- (-2.8369,-0.1228)
	-- (-2.8273,-0.4104)
	-- (-2.7845,-0.6942)
	-- (-2.7092,-0.9703)
	-- (-2.6025,-1.2351)
	-- (-2.4661,-1.4850)
	-- (-2.3020,-1.7168)
	-- (-2.1125,-1.9275)
	-- (-1.9003,-2.1141)
	-- (-1.6684,-2.2745)
	-- (-1.4199,-2.4065)
	-- (-1.1585,-2.5084)
	-- (-0.8875,-2.5791)
	-- (-0.6109,-2.6177)
	-- (-0.3323,-2.6239)
	-- (-0.0556,-2.5978)
	-- (0.2156,-2.5399)
	-- (0.4775,-2.4511)
	-- (0.7268,-2.3329)
	-- (0.9601,-2.1870)
	-- (1.1742,-2.0155)
	-- (1.3666,-1.8209)
	-- (1.5346,-1.6060)
	-- (1.6762,-1.3739)
	-- (1.7895,-1.1277)
	-- (1.8734,-0.8710)
	-- (1.9268,-0.6072)
	-- (1.9493,-0.3401)
	-- (1.9408,-0.0733)
	-- (1.9016,0.1897)
	-- (1.8326,0.4452)
	-- (1.7350,0.6898)
	-- (1.6102,0.9203)
	-- (1.4604,1.1338)
	-- (1.2878,1.3273)
	-- (1.0949,1.4985)
	-- (0.8847,1.6452)
	-- (0.6602,1.7657)
	-- (0.4248,1.8585)
	-- (0.1817,1.9228)
	-- (-0.0654,1.9579)
	-- (-0.3132,1.9637)
	-- (-0.5581,1.9404)
	-- (-0.7967,1.8888)
	-- (-1.0257,1.8099)
	-- (-1.2418,1.7053)
	-- (-1.4423,1.5767)
	-- (-1.6245,1.4263)
	-- (-1.7858,1.2567)
	-- (-1.9242,1.0704)
	-- (-2.0381,0.8706)
	-- (-2.1260,0.6603)
	-- (-2.1870,0.4427)
	-- (-2.2206,0.2214)
	-- (-2.2265,-0.0004)
	-- (-2.2050,-0.2193)
	-- (-2.1569,-0.4320)
	-- (-2.0831,-0.6352)
	-- (-1.9852,-0.8259)
	-- (-1.8648,-1.0013)
	-- (-1.7242,-1.1588)
	-- (-1.5658,-1.2960)
	-- (-1.3922,-1.4110)
	-- (-1.2064,-1.5021)
	-- (-1.0115,-1.5680)
	-- (-0.8106,-1.6078)
	-- (-0.6072,-1.6210)
	-- (-0.4044,-1.6075)
	-- (-0.2057,-1.5676)
	-- (-0.0143,-1.5021)
	-- (0.1667,-1.4120)
	-- (0.3342,-1.2989)
	-- (0.4854,-1.1646)
	-- (0.6179,-1.0112)
	-- (0.7293,-0.8412)
	-- (0.8177,-0.6574)
	-- (0.8815,-0.4627)
	-- (0.9194,-0.2603)
	-- (0.9307,-0.0533)
	-- (0.9148,0.1547)
	-- (0.8717,0.3606)
	-- (0.8018,0.5608)
	-- (0.7060,0.7521)
	-- (0.5853,0.9313)
	-- (0.4414,1.0954)
	-- (0.2762,1.2415)
	-- (0.0919,1.3670)
	-- (-0.1088,1.4696)
	-- (-0.3232,1.5474)
	-- (-0.5482,1.5987)
	-- (-0.7805,1.6223)
	-- (-1.0167,1.6173)
	-- (-1.2534,1.5833)
	-- (-1.4870,1.5202)
	-- (-1.7141,1.4285)
	-- (-1.9312,1.3090)
	-- (-2.1350,1.1628)
	-- (-2.3222,0.9917)
	-- (-2.4899,0.7976)
	-- (-2.6354,0.5828)
	-- (-2.7562,0.3500)
	-- (-2.8503,0.1023)
	-- (-2.9159,-0.1574)
	-- (-2.9515,-0.4255)
	-- (-2.9563,-0.6986)
	-- (-2.9296,-0.9730)
	-- (-2.8715,-1.2450)
	-- (-2.7821,-1.5110)
	-- (-2.6622,-1.7672)
	-- (-2.5131,-2.0102)
	-- (-2.3363,-2.2367)
	-- (-2.1338,-2.4433)
	-- (-1.9079,-2.6272)
	-- (-1.6613,-2.7856)
	-- (-1.3970,-2.9164)
	-- (-1.1183,-3.0174)
	-- (-0.8285,-3.0872)
	-- (-0.5315,-3.1245)
	-- (-0.2308,-3.1286)
	-- (0.0697,-3.0991)
	-- (0.3662,-3.0363)
	-- (0.6547,-2.9407)
	-- (0.9317,-2.8133)
	-- (1.1934,-2.6556)
	-- (1.4366,-2.4695)
	-- (1.6580,-2.2571)
	-- (1.8548,-2.0211)
	-- (2.0243,-1.7645)
	-- (2.1644,-1.4904)
	-- (2.2731,-1.2022)
	-- (2.3491,-0.9037)
	-- (2.3914,-0.5986)
	-- (2.3993,-0.2909)
	-- (2.3729,0.0157)
	-- (2.3123,0.3171)
	-- (2.2185,0.6095)
	-- (2.0927,0.8891)
	-- (1.9364,1.1524)
	-- (1.7519,1.3960)
	-- (1.5415,1.6169)
	-- (1.3079,1.8121)
	-- (1.0544,1.9794)
	-- (0.7842,2.1165)
	-- (0.5009,2.2217)
	-- (0.2082,2.2939)
	-- (-0.0899,2.3322)
	-- (-0.3897,2.3362)
	-- (-0.6871,2.3060)
	-- (-0.9782,2.2422)
	-- (-1.2592,2.1457)
	-- (-1.5265,2.0179)
	-- (-1.7766,1.8607)
	-- (-2.0063,1.6764)
	-- (-2.2125,1.4674)
	-- (-2.3926,1.2367)
	-- (-2.5445,0.9874)
	-- (-2.6661,0.7230)
	-- (-2.7560,0.4471)
	-- (-2.8132,0.1635)
	-- (-2.8370,-0.1241)
	-- (-2.8272,-0.4116)
	-- (-2.7842,-0.6954)
	-- (-2.7088,-0.9714)
	-- (-2.6020,-1.2362)
	-- (-2.4655,-1.4861)
	-- (-2.3013,-1.7178)
	-- (-2.1117,-1.9283)
	-- (-1.8994,-2.1149)
	-- (-1.6673,-2.2751)
	-- (-1.4188,-2.4069)
	-- (-1.1573,-2.5088)
	-- (-0.8864,-2.5793)
	-- (-0.6097,-2.6178)
	-- (-0.3312,-2.6239)
	-- (-0.0544,-2.5976)
	-- (0.2167,-2.5396)
	-- (0.4786,-2.4507)
	-- (0.7278,-2.3323)
	-- (0.9610,-2.1863)
	-- (1.1751,-2.0147)
	-- (1.3673,-1.8200)
	-- (1.5352,-1.6051)
	-- (1.6767,-1.3728)
	-- (1.7900,-1.1266)
	-- (1.8737,-0.8699)
	-- (1.9270,-0.6061)
	-- (1.9494,-0.3390)
	-- (1.9407,-0.0722)
	-- (1.9014,0.1908)
	-- (1.8323,0.4462)
	-- (1.7345,0.6908)
	-- (1.6097,0.9213)
	-- (1.4597,1.1346)
	-- (1.2870,1.3281)
	-- (1.0940,1.4992)
	-- (0.8838,1.6458)
	-- (0.6593,1.7661)
	-- (0.4238,1.8589)
	-- (0.1807,1.9230)
	-- (-0.0665,1.9580)
	-- (-0.3143,1.9636)
	-- (-0.5591,1.9403)
	-- (-0.7977,1.8885)
	-- (-1.0266,1.8096)
	-- (-1.2427,1.7048)
	-- (-1.4432,1.5761)
	-- (-1.6252,1.4257)
	-- (-1.7864,1.2559)
	-- (-1.9248,1.0696)
	-- (-2.0385,0.8697)
	-- (-2.1263,0.6594)
	-- (-2.1872,0.4418)
	-- (-2.2206,0.2204)
	-- (-2.2264,-0.0014)
	-- (-2.2049,-0.2202)
	-- (-2.1566,-0.4329)
	-- (-2.0828,-0.6360)
	-- (-1.9847,-0.8267)
	-- (-1.8643,-1.0020)
	-- (-1.7236,-1.1594)
	-- (-1.5651,-1.2966)
	-- (-1.3915,-1.4114)
	-- (-1.2056,-1.5024)
	-- (-1.0107,-1.5682)
	-- (-0.8098,-1.6079)
	-- (-0.6063,-1.6210)
	-- (-0.4035,-1.6074)
	-- (-0.2049,-1.5674)
	-- (-0.0135,-1.5018)
	-- (0.1674,-1.4116)
	-- (0.3349,-1.2984)
	-- (0.4860,-1.1640)
	-- (0.6184,-1.0105)
	-- (0.7297,-0.8405)
	-- (0.8180,-0.6566)
	-- (0.8817,-0.4619)
	-- (0.9195,-0.2594)
	-- (0.9306,-0.0525)
	-- (0.9146,0.1556)
	-- (0.8714,0.3615)
	-- (0.8015,0.5617)
	-- (0.7055,0.7529)
	-- (0.5848,0.9321)
	-- (0.4407,1.0960)
	-- (0.2754,1.2420)
	-- (0.0911,1.3674)
	-- (-0.1097,1.4700)
	-- (-0.3241,1.5476)
	-- (-0.5491,1.5988)
	-- (-0.7815,1.6223)
	-- (-1.0177,1.6172)
	-- (-1.2544,1.5831)
	-- (-1.4880,1.5199)
	-- (-1.7151,1.4281)
	-- (-1.9321,1.3084)
	-- (-2.1358,1.1621)
	-- (-2.3229,0.9909)
	-- (-2.4906,0.7967)
	-- (-2.6360,0.5818)
	-- (-2.7567,0.3490)
	-- (-2.8506,0.1012)
	-- (-2.9161,-0.1585)
	-- (-2.9516,-0.4267)
	-- (-2.9562,-0.6997)
	-- (-2.9295,-0.9741)
	-- (-2.8711,-1.2461)
	-- (-2.7816,-1.5121)
	-- (-2.6617,-1.7683)
	-- (-2.5124,-2.0112)
	-- (-2.3355,-2.2376)
	-- (-2.1329,-2.4441)
	-- (-1.9069,-2.6279)
	-- (-1.6602,-2.7862)
	-- (-1.3959,-2.9169)
	-- (-1.1171,-3.0178)
	-- (-0.8273,-3.0874)
	-- (-0.5302,-3.1246)
	-- (-0.2295,-3.1285)
	-- (0.0710,-3.0989)
	-- (0.3674,-3.0360)
	-- (0.6559,-2.9402)
	-- (0.9328,-2.8127)
	-- (1.1945,-2.6549)
	-- (1.4376,-2.4686)
	-- (1.6589,-2.2562)
	-- (1.8556,-2.0201)
	-- (2.0250,-1.7634)
	-- (2.1649,-1.4892)
	-- (2.2735,-1.2010)
	-- (2.3494,-0.9025)
	-- (2.3915,-0.5973)
	-- (2.3993,-0.2896)
	-- (2.3727,0.0170)
	-- (2.3120,0.3183)
	-- (2.2181,0.6107)
	-- (2.0921,0.8902)
	-- (1.9357,1.1535)
	-- (1.7510,1.3970)
	-- (1.5405,1.6178)
	-- (1.3069,1.8129)
	-- (1.0533,1.9800)
	-- (0.7830,2.1170)
	-- (0.4997,2.2221)
	-- (0.2070,2.2941)
	-- (-0.0912,2.3323)
	-- (-0.3909,2.3361)
	-- (-0.6883,2.3058)
	-- (-0.9794,2.2418)
	-- (-1.2604,2.1452)
	-- (-1.5276,2.0173)
	-- (-1.7776,1.8600)
	-- (-2.0072,1.6755)
	-- (-2.2133,1.4664)
	-- (-2.3933,1.2356)
	-- (-2.5451,0.9863)
	-- (-2.6666,0.7219)
	-- (-2.7563,0.4460)
	-- (-2.8133,0.1623)
	-- (-2.8370,-0.1253)
	-- (-2.8271,-0.4128)
	-- (-2.7840,-0.6965)
	-- (-2.7084,-0.9726)
	-- (-2.6015,-1.2373)
	-- (-2.4649,-1.4871)
	-- (-2.3005,-1.7187)
	-- (-2.1108,-1.9291)
	-- (-1.8984,-2.1156)
	-- (-1.6663,-2.2757)
	-- (-1.4178,-2.4074)
	-- (-1.1562,-2.5091)
	-- (-0.8852,-2.5795)
	-- (-0.6086,-2.6179)
	-- (-0.3300,-2.6238)
	-- (-0.0533,-2.5974)
	-- (0.2178,-2.5393)
	-- (0.4797,-2.4502)
	-- (0.7288,-2.3318)
	-- (0.9619,-2.1856)
	-- (1.1760,-2.0139)
	-- (1.3681,-1.8192)
	-- (1.5359,-1.6041)
	-- (1.6772,-1.3718)
	-- (1.7904,-1.1256)
	-- (1.8740,-0.8688)
	-- (1.9272,-0.6050)
	-- (1.9494,-0.3379)
	-- (1.9406,-0.0711)
	-- (1.9012,0.1918)
	-- (1.8319,0.4473)
	-- (1.7340,0.6918)
	-- (1.6091,0.9222)
	-- (1.4590,1.1355)
	-- (1.2862,1.3288)
	-- (1.0932,1.4998)
	-- (0.8829,1.6463)
	-- (0.6583,1.7666)
	-- (0.4228,1.8592)
	-- (0.1797,1.9232)
	-- (-0.0675,1.9581)
	-- (-0.3153,1.9636)
	-- (-0.5602,1.9401)
	-- (-0.7987,1.8883)
	-- (-1.0275,1.8092)
	-- (-1.2436,1.7043)
	-- (-1.4440,1.5755)
	-- (-1.6259,1.4250)
	-- (-1.7870,1.2552)
	-- (-1.9253,1.0688)
	-- (-2.0390,0.8688)
	-- (-2.1267,0.6584)
	-- (-2.1874,0.4409)
	-- (-2.2207,0.2195)
	-- (-2.2264,-0.0023)
	-- (-2.2047,-0.2212)
	-- (-2.1564,-0.4337)
	-- (-2.0824,-0.6369)
	-- (-1.9842,-0.8275)
	-- (-1.8637,-1.0027)
	-- (-1.7230,-1.1600)
	-- (-1.5644,-1.2971)
	-- (-1.3907,-1.4119)
	-- (-1.2048,-1.5028)
	-- (-1.0098,-1.5684)
	-- (-0.8089,-1.6080)
	-- (-0.6054,-1.6210)
	-- (-0.4027,-1.6073)
	-- (-0.2040,-1.5672)
	-- (-0.0127,-1.5015)
	-- (0.1682,-1.4112)
	-- (0.3355,-1.2979)
	-- (0.4866,-1.1633)
	-- (0.6189,-1.0098)
	-- (0.7302,-0.8397)
	-- (0.8183,-0.6558)
	-- (0.8819,-0.4610)
	-- (0.9196,-0.2585)
	-- (0.9306,-0.0516)
	-- (0.9145,0.1565)
	-- (0.8712,0.3623)
	-- (0.8011,0.5625)
	-- (0.7051,0.7537)
	-- (0.5842,0.9328)
	-- (0.4401,1.0967)
	-- (0.2747,1.2426)
	-- (0.0903,1.3679)
	-- (-0.1106,1.4704)
	-- (-0.3251,1.5479)
	-- (-0.5501,1.5990)
	-- (-0.7824,1.6224)
	-- (-1.0187,1.6171)
	-- (-1.2554,1.5829)
	-- (-1.4890,1.5196)
	-- (-1.7160,1.4276)
	-- (-1.9330,1.3078)
	-- (-2.1366,1.1615)
	-- (-2.3237,0.9901)
	-- (-2.4912,0.7958)
	-- (-2.6365,0.5809)
	-- (-2.7571,0.3480)
	-- (-2.8510,0.1001)
	-- (-2.9163,-0.1596)
	-- (-2.9517,-0.4278)
	-- (-2.9562,-0.7009)
	-- (-2.9293,-0.9753)
	-- (-2.8708,-1.2473)
	-- (-2.7812,-1.5132)
	-- (-2.6611,-1.7693)
	-- (-2.5117,-2.0122)
	-- (-2.3347,-2.2385)
	-- (-2.1320,-2.4449)
	-- (-1.9059,-2.6286)
	-- (-1.6591,-2.7868)
	-- (-1.3947,-2.9174)
	-- (-1.1159,-3.0181)
	-- (-0.8261,-3.0876)
	-- (-0.5289,-3.1247)
	-- (-0.2282,-3.1285)
	-- (0.0723,-3.0988)
	-- (0.3686,-3.0357)
	-- (0.6571,-2.9398)
	-- (0.9339,-2.8121)
	-- (1.1956,-2.6542)
	-- (1.4386,-2.4678)
	-- (1.6598,-2.2552)
	-- (1.8563,-2.0191)
	-- (2.0256,-1.7622)
	-- (2.1654,-1.4880)
	-- (2.2739,-1.1998)
	-- (2.3496,-0.9012)
	-- (2.3916,-0.5961)
	-- (2.3993,-0.2883)
	-- (2.3725,0.0183)
	-- (2.3117,0.3196)
	-- (2.2176,0.6119)
	-- (2.0915,0.8914)
	-- (1.9350,1.1546)
	-- (1.7502,1.3980)
	-- (1.5396,1.6186)
	-- (1.3059,1.8137)
	-- (1.0522,1.9806)
	-- (0.7819,2.1175)
	-- (0.4985,2.2225)
	-- (0.2057,2.2944)
	-- (-0.0925,2.3324)
	-- (-0.3922,2.3361)
	-- (-0.6895,2.3056)
	-- (-0.9806,2.2415)
	-- (-1.2615,2.1447)
	-- (-1.5287,2.0167)
	-- (-1.7787,1.8593)
	-- (-2.0081,1.6747)
	-- (-2.2141,1.4655)
	-- (-2.3940,1.2346)
	-- (-2.5456,0.9852)
	-- (-2.6670,0.7207)
	-- (-2.7566,0.4448)
	-- (-2.8135,0.1611)
	-- (-2.8370,-0.1265)
	-- (-2.8270,-0.4140)
	-- (-2.7837,-0.6977)
	-- (-2.7080,-0.9737)
	-- (-2.6010,-1.2383)
	-- (-2.4642,-1.4881)
	-- (-2.2998,-1.7196)
	-- (-2.1100,-1.9300)
	-- (-1.8975,-2.1163)
	-- (-1.6653,-2.2763)
	-- (-1.4167,-2.4079)
	-- (-1.1551,-2.5095)
	-- (-0.8841,-2.5798)
	-- (-0.6074,-2.6180)
	-- (-0.3288,-2.6238)
	-- (-0.0521,-2.5973)
	-- (0.2190,-2.5389)
	-- (0.4808,-2.4498)
	-- (0.7299,-2.3312)
	-- (0.9629,-2.1849)
	-- (1.1768,-2.0131)
	-- (1.3689,-1.8183)
	-- (1.5365,-1.6032)
	-- (1.6778,-1.3708)
	-- (1.7908,-1.1245)
	-- (1.8743,-0.8677)
	-- (1.9273,-0.6039)
	-- (1.9494,-0.3368)
	-- (1.9405,-0.0699)
	-- (1.9010,0.1929)
	-- (1.8316,0.4483)
	-- (1.7336,0.6928)
	-- (1.6085,0.9232)
	-- (1.4584,1.1363)
	-- (1.2854,1.3296)
	-- (1.0923,1.5005)
	-- (0.8819,1.6469)
	-- (0.6573,1.7670)
	-- (0.4217,1.8595)
	-- (0.1786,1.9234)
	-- (-0.0686,1.9581)
	-- (-0.3163,1.9636)
	-- (-0.5612,1.9399)
	-- (-0.7997,1.8880)
	-- (-1.0285,1.8088)
	-- (-1.2445,1.7038)
	-- (-1.4448,1.5749)
	-- (-1.6266,1.4243)
	-- (-1.7877,1.2544)
	-- (-1.9258,1.0680)
	-- (-2.0394,0.8680)
	-- (-2.1270,0.6575)
	-- (-2.1876,0.4400)
	-- (-2.2208,0.2186)
	-- (-2.2264,-0.0032)
	-- (-2.2046,-0.2221)
	-- (-2.1561,-0.4346)
	-- (-2.0820,-0.6377)
	-- (-1.9838,-0.8282)
	-- (-1.8632,-1.0034)
	-- (-1.7223,-1.1607)
	-- (-1.5637,-1.2976)
	-- (-1.3900,-1.4123)
	-- (-1.2040,-1.5031)
	-- (-1.0090,-1.5686)
	-- (-0.8081,-1.6081)
	-- (-0.6046,-1.6210)
	-- (-0.4019,-1.6072)
	-- (-0.2032,-1.5670)
	-- (-0.0119,-1.5011)
	-- (0.1689,-1.4107)
	-- (0.3362,-1.2973)
	-- (0.4872,-1.1627)
	-- (0.6195,-1.0091)
	-- (0.7306,-0.8390)
	-- (0.8187,-0.6550)
	-- (0.8821,-0.4602)
	-- (0.9197,-0.2577)
	-- (0.9306,-0.0507)
	-- (0.9144,0.1574)
	-- (0.8710,0.3632)
	-- (0.8008,0.5633)
	-- (0.7046,0.7545)
	-- (0.5836,0.9335)
	-- (0.4394,1.0974)
	-- (0.2740,1.2432)
	-- (0.0895,1.3684)
	-- (-0.1115,1.4707)
	-- (-0.3260,1.5482)
	-- (-0.5511,1.5992)
	-- (-0.7834,1.6224)
	-- (-1.0197,1.6170)
	-- (-1.2564,1.5827)
	-- (-1.4900,1.5192)
	-- (-1.7169,1.4272)
	-- (-1.9339,1.3073)
	-- (-2.1374,1.1608)
	-- (-2.3244,0.9894)
	-- (-2.4919,0.7950)
	-- (-2.6371,0.5800)
	-- (-2.7576,0.3470)
	-- (-2.8513,0.0990)
	-- (-2.9165,-0.1607)
	-- (-2.9518,-0.4289)
	-- (-2.9561,-0.7021)
	-- (-2.9291,-0.9764)
	-- (-2.8705,-1.2484)
	-- (-2.7808,-1.5143)
	-- (-2.6605,-1.7704)
	-- (-2.5110,-2.0132)
	-- (-2.3339,-2.2394)
	-- (-2.1310,-2.4457)
	-- (-1.9049,-2.6293)
	-- (-1.6581,-2.7875)
	-- (-1.3936,-2.9178)
	-- (-1.1147,-3.0185)
	-- (-0.8248,-3.0879)
	-- (-0.5277,-3.1247)
	-- (-0.2269,-3.1284)
	-- (0.0735,-3.0986)
	-- (0.3699,-3.0353)
	-- (0.6583,-2.9393)
	-- (0.9351,-2.8115)
	-- (1.1966,-2.6535)
	-- (1.4396,-2.4669)
	-- (1.6607,-2.2543)
	-- (1.8571,-2.0180)
	-- (2.0263,-1.7611)
	-- (2.1659,-1.4868)
	-- (2.2743,-1.1985)
	-- (2.3499,-0.8999)
	-- (2.3917,-0.5948)
	-- (2.3992,-0.2870)
	-- (2.3723,0.0195)
	-- (2.3114,0.3208)
	-- (2.2171,0.6131)
	-- (2.0909,0.8925)
	-- (1.9343,1.1556)
	-- (1.7494,1.3990)
	-- (1.5386,1.6195)
	-- (1.3048,1.8144)
	-- (1.0511,1.9813)
	-- (0.7807,2.1180)
	-- (0.4972,2.2228)
	-- (0.2045,2.2946)
	-- (-0.0937,2.3325)
	-- (-0.3935,2.3360)
	-- (-0.6908,2.3054)
	-- (-0.9818,2.2412)
	-- (-1.2627,2.1443)
	-- (-1.5298,2.0161)
	-- (-1.7797,1.8586)
	-- (-2.0090,1.6739)
	-- (-2.2149,1.4646)
	-- (-2.3947,1.2336)
	-- (-2.5462,0.9842)
	-- (-2.6674,0.7196)
	-- (-2.7569,0.4436)
	-- (-2.8137,0.1599)
	-- (-2.8370,-0.1277)
	-- (-2.8269,-0.4152)
	-- (-2.7835,-0.6989)
	-- (-2.7076,-0.9748)
	-- (-2.6005,-1.2394)
	-- (-2.4636,-1.4891)
	-- (-2.2990,-1.7206)
	-- (-2.1091,-1.9308)
	-- (-1.8965,-2.1171)
	-- (-1.6643,-2.2770)
	-- (-1.4156,-2.4084)
	-- (-1.1540,-2.5098)
	-- (-0.8829,-2.5800)
	-- (-0.6062,-2.6181)
	-- (-0.3277,-2.6237)
	-- (-0.0510,-2.5971)
	-- (0.2201,-2.5386)
	-- (0.4819,-2.4494)
	-- (0.7309,-2.3306)
	-- (0.9638,-2.1843)
	-- (1.1777,-2.0124)
	-- (1.3696,-1.8174)
	-- (1.5372,-1.6022)
	-- (1.6783,-1.3698)
	-- (1.7912,-1.1235)
	-- (1.8746,-0.8666)
	-- (1.9275,-0.6028)
	-- (1.9494,-0.3356)
	-- (1.9404,-0.0688)
	-- (1.9007,0.1940)
	-- (1.8312,0.4494)
	-- (1.7331,0.6938)
	-- (1.6079,0.9241)
	-- (1.4577,1.1372)
	-- (1.2847,1.3304)
	-- (1.0915,1.5012)
	-- (0.8810,1.6475)
	-- (0.6563,1.7675)
	-- (0.4207,1.8599)
	-- (0.1776,1.9236)
	-- (-0.0696,1.9582)
	-- (-0.3174,1.9635)
	-- (-0.5622,1.9398)
	-- (-0.8006,1.8877)
	-- (-1.0294,1.8084)
	-- (-1.2454,1.7033)
	-- (-1.4456,1.5743)
	-- (-1.6273,1.4236)
	-- (-1.7883,1.2537)
	-- (-1.9264,1.0672)
	-- (-2.0398,0.8671)
	-- (-2.1273,0.6566)
	-- (-2.1878,0.4390)
	-- (-2.2209,0.2176)
	-- (-2.2263,-0.0042)
	-- (-2.2044,-0.2230)
	-- (-2.1558,-0.4355)
	-- (-2.0817,-0.6385)
	-- (-1.9833,-0.8290)
	-- (-1.8626,-1.0041)
	-- (-1.7217,-1.1613)
	-- (-1.5630,-1.2982)
	-- (-1.3892,-1.4127)
	-- (-1.2032,-1.5034)
	-- (-1.0082,-1.5689)
	-- (-0.8072,-1.6082)
	-- (-0.6037,-1.6210)
	-- (-0.4010,-1.6070)
	-- (-0.2024,-1.5667)
	-- (-0.0111,-1.5008)
	-- (0.1696,-1.4103)
	-- (0.3369,-1.2968)
	-- (0.4878,-1.1621)
	-- (0.6200,-1.0084)
	-- (0.7310,-0.8382)
	-- (0.8190,-0.6542)
	-- (0.8823,-0.4594)
	-- (0.9198,-0.2568)
	-- (0.9306,-0.0498)
	-- (0.9143,0.1582)
	-- (0.8707,0.3640)
	-- (0.8004,0.5641)
	-- (0.7041,0.7553)
	-- (0.5831,0.9342)
	-- (0.4388,1.0980)
	-- (0.2732,1.2437)
	-- (0.0887,1.3689)
	-- (-0.1123,1.4711)
	-- (-0.3269,1.5485)
	-- (-0.5520,1.5993)
	-- (-0.7844,1.6224)
	-- (-1.0207,1.6170)
	-- (-1.2574,1.5825)
	-- (-1.4909,1.5189)
	-- (-1.7179,1.4267)
	-- (-1.9348,1.3067)
	-- (-2.1383,1.1601)
	-- (-2.3252,0.9886)
	-- (-2.4925,0.7941)
	-- (-2.6376,0.5790)
	-- (-2.7581,0.3460)
	-- (-2.8517,0.0980)
	-- (-2.9167,-0.1618)
	-- (-2.9518,-0.4301)
	-- (-2.9561,-0.7032)
	-- (-2.9289,-0.9776)
	-- (-2.8702,-1.2495)
	-- (-2.7803,-1.5154)
	-- (-2.6600,-1.7714)
	-- (-2.5104,-2.0142)
	-- (-2.3331,-2.2403)
	-- (-2.1301,-2.4466)
	-- (-1.9039,-2.6300)
	-- (-1.6570,-2.7881)
	-- (-1.3924,-2.9183)
	-- (-1.1135,-3.0189)
	-- (-0.8236,-3.0881)
	-- (-0.5264,-3.1248)
	-- (-0.2257,-3.1284)
	-- (0.0748,-3.0984)
	-- (0.3711,-3.0350)
	-- (0.6595,-2.9388)
	-- (0.9362,-2.8109)
	-- (1.1977,-2.6527)
	-- (1.4405,-2.4661)
}
\def\scaleC{0.4620}
\def\bildkurveD{
	(0.4392,0.4048)
	-- (0.3640,0.5001)
	-- (0.2780,0.5918)
	-- (0.1813,0.6792)
	-- (0.0740,0.7615)
	-- (-0.0437,0.8379)
	-- (-0.1715,0.9076)
	-- (-0.3091,0.9700)
	-- (-0.4562,1.0242)
	-- (-0.6124,1.0696)
	-- (-0.7774,1.1057)
	-- (-0.9507,1.1319)
	-- (-1.1318,1.1476)
	-- (-1.3202,1.1526)
	-- (-1.5157,1.1463)
	-- (-1.7176,1.1286)
	-- (-1.9255,1.0992)
	-- (-2.1391,1.0582)
	-- (-2.3580,1.0055)
	-- (-2.5817,0.9413)
	-- (-2.8102,0.8659)
	-- (-3.0431,0.7800)
	-- (-3.2803,0.6841)
	-- (-3.5217,0.5793)
	-- (-3.7673,0.4668)
	-- (-4.0172,0.3483)
	-- (-4.2713,0.2258)
	-- (-4.5295,0.1018)
	-- (-4.7917,-0.0205)
	-- (-5.0571,-0.1375)
	-- (-5.3246,-0.2447)
	-- (-5.5922,-0.3372)
	-- (-5.8563,-0.4095)
	-- (-6.1120,-0.4560)
	-- (-6.3522,-0.4717)
	-- (-6.5675,-0.4532)
	-- (-6.7470,-0.3994)
	-- (-6.8785,-0.3130)
	-- (-6.9507,-0.2015)
	-- (-6.9549,-0.0762)
	-- (-6.8874,0.0480)
	-- (-6.7504,0.1561)
	-- (-6.5517,0.2345)
	-- (-6.3031,0.2742)
	-- (-6.0186,0.2708)
	-- (-5.7114,0.2252)
	-- (-5.3930,0.1415)
	-- (-5.0721,0.0262)
	-- (-4.7545,-0.1132)
	-- (-4.4438,-0.2694)
	-- (-4.1418,-0.4359)
	-- (-3.8491,-0.6069)
	-- (-3.5655,-0.7777)
	-- (-3.2903,-0.9445)
	-- (-3.0229,-1.1045)
	-- (-2.7624,-1.2551)
	-- (-2.5082,-1.3949)
	-- (-2.2598,-1.5224)
	-- (-2.0167,-1.6367)
	-- (-1.7789,-1.7371)
	-- (-1.5461,-1.8234)
	-- (-1.3184,-1.8951)
	-- (-1.0961,-1.9523)
	-- (-0.8793,-1.9950)
	-- (-0.6685,-2.0234)
	-- (-0.4639,-2.0377)
	-- (-0.2660,-2.0384)
	-- (-0.0754,-2.0257)
	-- (0.1076,-2.0002)
	-- (0.2824,-1.9626)
	-- (0.4485,-1.9133)
	-- (0.6055,-1.8530)
	-- (0.7528,-1.7826)
	-- (0.8902,-1.7026)
	-- (1.0172,-1.6139)
	-- (1.1334,-1.5174)
	-- (1.2386,-1.4139)
	-- (1.3324,-1.3042)
	-- (1.4148,-1.1892)
	-- (1.4855,-1.0699)
	-- (1.5445,-0.9471)
	-- (1.5917,-0.8218)
	-- (1.6271,-0.6949)
	-- (1.6510,-0.5672)
	-- (1.6633,-0.4396)
	-- (1.6644,-0.3131)
	-- (1.6545,-0.1885)
	-- (1.6340,-0.0666)
	-- (1.6031,0.0519)
	-- (1.5625,0.1661)
	-- (1.5125,0.2753)
	-- (1.4537,0.3789)
	-- (1.3868,0.4763)
	-- (1.3123,0.5669)
	-- (1.2308,0.6501)
	-- (1.1432,0.7256)
	-- (1.0501,0.7928)
	-- (0.9522,0.8516)
	-- (0.8504,0.9015)
	-- (0.7455,0.9423)
	-- (0.6381,0.9740)
	-- (0.5292,0.9965)
	-- (0.4195,1.0096)
	-- (0.3098,1.0136)
	-- (0.2009,1.0083)
	-- (0.0935,0.9942)
	-- (-0.0116,0.9713)
	-- (-0.1136,0.9400)
	-- (-0.2120,0.9005)
	-- (-0.3060,0.8534)
	-- (-0.3950,0.7991)
	-- (-0.4785,0.7380)
	-- (-0.5559,0.6708)
	-- (-0.6268,0.5979)
	-- (-0.6907,0.5201)
	-- (-0.7472,0.4379)
	-- (-0.7961,0.3521)
	-- (-0.8371,0.2633)
	-- (-0.8700,0.1723)
	-- (-0.8947,0.0797)
	-- (-0.9111,-0.0136)
	-- (-0.9192,-0.1069)
	-- (-0.9190,-0.1995)
	-- (-0.9107,-0.2907)
	-- (-0.8944,-0.3797)
	-- (-0.8704,-0.4658)
	-- (-0.8390,-0.5484)
	-- (-0.8006,-0.6269)
	-- (-0.7555,-0.7005)
	-- (-0.7042,-0.7688)
	-- (-0.6473,-0.8311)
	-- (-0.5852,-0.8870)
	-- (-0.5187,-0.9361)
	-- (-0.4482,-0.9778)
	-- (-0.3745,-1.0120)
	-- (-0.2983,-1.0383)
	-- (-0.2203,-1.0564)
	-- (-0.1412,-1.0662)
	-- (-0.0618,-1.0676)
	-- (0.0171,-1.0606)
	-- (0.0949,-1.0451)
	-- (0.1706,-1.0212)
	-- (0.2435,-0.9891)
	-- (0.3130,-0.9490)
	-- (0.3782,-0.9011)
	-- (0.4384,-0.8458)
	-- (0.4930,-0.7834)
	-- (0.5412,-0.7144)
	-- (0.5825,-0.6394)
	-- (0.6162,-0.5587)
	-- (0.6418,-0.4731)
	-- (0.6587,-0.3831)
	-- (0.6666,-0.2895)
	-- (0.6649,-0.1929)
	-- (0.6534,-0.0941)
	-- (0.6317,0.0061)
	-- (0.5995,0.1070)
	-- (0.5568,0.2077)
	-- (0.5032,0.3074)
	-- (0.4389,0.4052)
	-- (0.3636,0.5004)
	-- (0.2776,0.5922)
	-- (0.1808,0.6796)
	-- (0.0735,0.7619)
	-- (-0.0442,0.8382)
	-- (-0.1720,0.9079)
	-- (-0.3097,0.9702)
	-- (-0.4568,1.0244)
	-- (-0.6131,1.0698)
	-- (-0.7781,1.1058)
	-- (-0.9514,1.1320)
	-- (-1.1325,1.1477)
	-- (-1.3211,1.1526)
	-- (-1.5165,1.1463)
	-- (-1.7184,1.1285)
	-- (-1.9264,1.0991)
	-- (-2.1400,1.0580)
	-- (-2.3589,1.0052)
	-- (-2.5827,0.9410)
	-- (-2.8111,0.8656)
	-- (-3.0441,0.7796)
	-- (-3.2813,0.6837)
	-- (-3.5227,0.5788)
	-- (-3.7684,0.4663)
	-- (-4.0182,0.3478)
	-- (-4.2723,0.2253)
	-- (-4.5306,0.1013)
	-- (-4.7928,-0.0210)
	-- (-5.0582,-0.1380)
	-- (-5.3258,-0.2451)
	-- (-5.5933,-0.3375)
	-- (-5.8574,-0.4097)
	-- (-6.1130,-0.4561)
	-- (-6.3531,-0.4717)
	-- (-6.5684,-0.4530)
	-- (-6.7477,-0.3991)
	-- (-6.8789,-0.3126)
	-- (-6.9508,-0.2010)
	-- (-6.9548,-0.0757)
	-- (-6.8870,0.0485)
	-- (-6.7497,0.1565)
	-- (-6.5507,0.2348)
	-- (-6.3020,0.2743)
	-- (-6.0173,0.2707)
	-- (-5.7101,0.2249)
	-- (-5.3917,0.1410)
	-- (-5.0707,0.0257)
	-- (-4.7532,-0.1138)
	-- (-4.4425,-0.2701)
	-- (-4.1406,-0.4366)
	-- (-3.8479,-0.6076)
	-- (-3.5643,-0.7784)
	-- (-3.2892,-0.9452)
	-- (-3.0218,-1.1051)
	-- (-2.7613,-1.2558)
	-- (-2.5071,-1.3954)
	-- (-2.2587,-1.5229)
	-- (-2.0157,-1.6371)
	-- (-1.7779,-1.7375)
	-- (-1.5451,-1.8237)
	-- (-1.3175,-1.8954)
	-- (-1.0952,-1.9525)
	-- (-0.8784,-1.9952)
	-- (-0.6676,-2.0235)
	-- (-0.4630,-2.0378)
	-- (-0.2652,-2.0383)
	-- (-0.0746,-2.0256)
	-- (0.1083,-2.0001)
	-- (0.2831,-1.9624)
	-- (0.4492,-1.9130)
	-- (0.6061,-1.8528)
	-- (0.7534,-1.7822)
	-- (0.8908,-1.7022)
	-- (1.0177,-1.6136)
	-- (1.1339,-1.5170)
	-- (1.2390,-1.4134)
	-- (1.3328,-1.3037)
	-- (1.4151,-1.1887)
	-- (1.4858,-1.0694)
	-- (1.5447,-0.9466)
	-- (1.5918,-0.8213)
	-- (1.6273,-0.6943)
	-- (1.6510,-0.5666)
	-- (1.6634,-0.4391)
	-- (1.6644,-0.3126)
	-- (1.6544,-0.1880)
	-- (1.6338,-0.0661)
	-- (1.6030,0.0523)
	-- (1.5623,0.1665)
	-- (1.5123,0.2758)
	-- (1.4535,0.3794)
	-- (1.3865,0.4767)
	-- (1.3119,0.5673)
	-- (1.2305,0.6505)
	-- (1.1428,0.7259)
	-- (1.0497,0.7931)
	-- (0.9518,0.8518)
	-- (0.8500,0.9017)
	-- (0.7450,0.9425)
	-- (0.6377,0.9742)
	-- (0.5287,0.9966)
	-- (0.4190,1.0097)
	-- (0.3093,1.0136)
	-- (0.2004,1.0083)
	-- (0.0931,0.9941)
	-- (-0.0120,0.9712)
	-- (-0.1140,0.9398)
	-- (-0.2124,0.9004)
	-- (-0.3063,0.8532)
	-- (-0.3954,0.7989)
	-- (-0.4788,0.7378)
	-- (-0.5562,0.6705)
	-- (-0.6271,0.5976)
	-- (-0.6909,0.5197)
	-- (-0.7475,0.4375)
	-- (-0.7963,0.3517)
	-- (-0.8373,0.2629)
	-- (-0.8702,0.1719)
	-- (-0.8948,0.0793)
	-- (-0.9111,-0.0140)
	-- (-0.9192,-0.1073)
	-- (-0.9190,-0.1999)
	-- (-0.9106,-0.2910)
	-- (-0.8943,-0.3800)
	-- (-0.8703,-0.4662)
	-- (-0.8389,-0.5488)
	-- (-0.8004,-0.6272)
	-- (-0.7553,-0.7008)
	-- (-0.7040,-0.7690)
	-- (-0.6471,-0.8313)
	-- (-0.5850,-0.8872)
	-- (-0.5184,-0.9362)
	-- (-0.4479,-0.9780)
	-- (-0.3742,-1.0121)
	-- (-0.2980,-1.0384)
	-- (-0.2199,-1.0564)
	-- (-0.1409,-1.0662)
	-- (-0.0615,-1.0676)
	-- (0.0175,-1.0605)
	-- (0.0952,-1.0450)
	-- (0.1709,-1.0211)
	-- (0.2438,-0.9889)
	-- (0.3133,-0.9488)
	-- (0.3785,-0.9009)
	-- (0.4387,-0.8455)
	-- (0.4932,-0.7831)
	-- (0.5414,-0.7141)
	-- (0.5827,-0.6390)
	-- (0.6163,-0.5584)
	-- (0.6419,-0.4727)
	-- (0.6588,-0.3827)
	-- (0.6666,-0.2891)
	-- (0.6649,-0.1925)
	-- (0.6533,-0.0937)
	-- (0.6316,0.0065)
	-- (0.5994,0.1074)
	-- (0.5565,0.2081)
	-- (0.5030,0.3078)
	-- (0.4386,0.4056)
	-- (0.3633,0.5008)
	-- (0.2772,0.5925)
	-- (0.1804,0.6799)
	-- (0.0730,0.7622)
	-- (-0.0447,0.8385)
	-- (-0.1726,0.9082)
	-- (-0.3103,0.9704)
	-- (-0.4575,1.0246)
	-- (-0.6138,1.0700)
	-- (-0.7788,1.1060)
	-- (-0.9521,1.1321)
	-- (-1.1333,1.1477)
	-- (-1.3219,1.1526)
	-- (-1.5173,1.1462)
	-- (-1.7193,1.1284)
	-- (-1.9273,1.0989)
	-- (-2.1409,1.0578)
	-- (-2.3598,1.0050)
	-- (-2.5836,0.9407)
	-- (-2.8121,0.8652)
	-- (-3.0450,0.7792)
	-- (-3.2823,0.6832)
	-- (-3.5237,0.5784)
	-- (-3.7694,0.4658)
	-- (-4.0193,0.3473)
	-- (-4.2734,0.2248)
	-- (-4.5317,0.1008)
	-- (-4.7939,-0.0215)
	-- (-5.0593,-0.1384)
	-- (-5.3269,-0.2456)
	-- (-5.5944,-0.3379)
	-- (-5.8585,-0.4100)
	-- (-6.1141,-0.4562)
	-- (-6.3541,-0.4717)
	-- (-6.5692,-0.4529)
	-- (-6.7483,-0.3988)
	-- (-6.8794,-0.3122)
	-- (-6.9510,-0.2005)
	-- (-6.9546,-0.0752)
	-- (-6.8866,0.0490)
	-- (-6.7490,0.1569)
	-- (-6.5498,0.2350)
	-- (-6.3008,0.2743)
	-- (-6.0161,0.2706)
	-- (-5.7088,0.2246)
	-- (-5.3903,0.1406)
	-- (-5.0694,0.0251)
	-- (-4.7518,-0.1144)
	-- (-4.4412,-0.2708)
	-- (-4.1393,-0.4373)
	-- (-3.8467,-0.6083)
	-- (-3.5631,-0.7791)
	-- (-3.2880,-0.9459)
	-- (-3.0206,-1.1058)
	-- (-2.7602,-1.2564)
	-- (-2.5061,-1.3960)
	-- (-2.2577,-1.5234)
	-- (-2.0147,-1.6376)
	-- (-1.7769,-1.7379)
	-- (-1.5441,-1.8240)
	-- (-1.3165,-1.8957)
	-- (-1.0943,-1.9527)
	-- (-0.8775,-1.9953)
	-- (-0.6667,-2.0236)
	-- (-0.4622,-2.0378)
	-- (-0.2644,-2.0383)
	-- (-0.0738,-2.0255)
	-- (0.1091,-2.0000)
	-- (0.2838,-1.9622)
	-- (0.4498,-1.9128)
	-- (0.6067,-1.8525)
	-- (0.7540,-1.7819)
	-- (0.8913,-1.7019)
	-- (1.0182,-1.6132)
	-- (1.1343,-1.5166)
	-- (1.2394,-1.4130)
	-- (1.3332,-1.3032)
	-- (1.4154,-1.1882)
	-- (1.4860,-1.0689)
	-- (1.5449,-0.9461)
	-- (1.5920,-0.8207)
	-- (1.6274,-0.6938)
	-- (1.6511,-0.5661)
	-- (1.6634,-0.4386)
	-- (1.6644,-0.3121)
	-- (1.6544,-0.1875)
	-- (1.6337,-0.0656)
	-- (1.6028,0.0528)
	-- (1.5621,0.1670)
	-- (1.5120,0.2762)
	-- (1.4532,0.3798)
	-- (1.3862,0.4771)
	-- (1.3116,0.5676)
	-- (1.2301,0.6508)
	-- (1.1424,0.7262)
	-- (1.0493,0.7934)
	-- (0.9514,0.8520)
	-- (0.8496,0.9018)
	-- (0.7446,0.9426)
	-- (0.6372,0.9743)
	-- (0.5283,0.9966)
	-- (0.4186,1.0097)
	-- (0.3089,1.0135)
	-- (0.2000,1.0083)
	-- (0.0926,0.9940)
	-- (-0.0124,0.9711)
	-- (-0.1145,0.9397)
	-- (-0.2128,0.9002)
	-- (-0.3067,0.8530)
	-- (-0.3957,0.7986)
	-- (-0.4792,0.7375)
	-- (-0.5565,0.6702)
	-- (-0.6273,0.5973)
	-- (-0.6912,0.5194)
	-- (-0.7477,0.4372)
	-- (-0.7965,0.3513)
	-- (-0.8375,0.2625)
	-- (-0.8703,0.1715)
	-- (-0.8949,0.0789)
	-- (-0.9112,-0.0143)
	-- (-0.9192,-0.1077)
	-- (-0.9189,-0.2003)
	-- (-0.9106,-0.2914)
	-- (-0.8942,-0.3804)
	-- (-0.8702,-0.4665)
	-- (-0.8387,-0.5491)
	-- (-0.8002,-0.6275)
	-- (-0.7551,-0.7011)
	-- (-0.7038,-0.7693)
	-- (-0.6468,-0.8316)
	-- (-0.5847,-0.8874)
	-- (-0.5181,-0.9364)
	-- (-0.4476,-0.9782)
	-- (-0.3739,-1.0123)
	-- (-0.2976,-1.0384)
	-- (-0.2196,-1.0565)
	-- (-0.1405,-1.0663)
	-- (-0.0611,-1.0676)
	-- (0.0178,-1.0605)
	-- (0.0955,-1.0449)
	-- (0.1712,-1.0210)
	-- (0.2441,-0.9888)
	-- (0.3136,-0.9486)
	-- (0.3787,-0.9007)
	-- (0.4389,-0.8453)
	-- (0.4934,-0.7829)
	-- (0.5416,-0.7138)
	-- (0.5828,-0.6387)
	-- (0.6165,-0.5580)
	-- (0.6420,-0.4724)
	-- (0.6588,-0.3824)
	-- (0.6666,-0.2887)
	-- (0.6649,-0.1921)
	-- (0.6532,-0.0933)
	-- (0.6314,0.0069)
	-- (0.5992,0.1078)
	-- (0.5563,0.2085)
	-- (0.5027,0.3082)
	-- (0.4383,0.4061)
	-- (0.3630,0.5012)
	-- (0.2768,0.5929)
	-- (0.1800,0.6803)
	-- (0.0726,0.7625)
	-- (-0.0452,0.8388)
	-- (-0.1732,0.9085)
	-- (-0.3109,0.9707)
	-- (-0.4581,1.0248)
	-- (-0.6145,1.0701)
	-- (-0.7795,1.1061)
	-- (-0.9529,1.1322)
	-- (-1.1341,1.1478)
	-- (-1.3227,1.1526)
	-- (-1.5182,1.1462)
	-- (-1.7202,1.1283)
	-- (-1.9282,1.0988)
	-- (-2.1418,1.0576)
	-- (-2.3608,1.0047)
	-- (-2.5846,0.9404)
	-- (-2.8131,0.8649)
	-- (-3.0460,0.7788)
	-- (-3.2833,0.6828)
	-- (-3.5248,0.5779)
	-- (-3.7704,0.4654)
	-- (-4.0203,0.3468)
	-- (-4.2745,0.2242)
	-- (-4.5328,0.1003)
	-- (-4.7950,-0.0220)
	-- (-5.0605,-0.1389)
	-- (-5.3280,-0.2460)
	-- (-5.5955,-0.3382)
	-- (-5.8596,-0.4102)
	-- (-6.1151,-0.4564)
	-- (-6.3551,-0.4717)
	-- (-6.5701,-0.4527)
	-- (-6.7490,-0.3985)
	-- (-6.8798,-0.3118)
	-- (-6.9512,-0.1999)
	-- (-6.9545,-0.0746)
	-- (-6.8861,0.0495)
	-- (-6.7483,0.1573)
	-- (-6.5488,0.2353)
	-- (-6.2997,0.2744)
	-- (-6.0148,0.2705)
	-- (-5.7074,0.2243)
	-- (-5.3890,0.1402)
	-- (-5.0680,0.0246)
	-- (-4.7505,-0.1151)
	-- (-4.4399,-0.2715)
	-- (-4.1381,-0.4380)
	-- (-3.8455,-0.6090)
	-- (-3.5620,-0.7798)
	-- (-3.2869,-0.9466)
	-- (-3.0195,-1.1064)
	-- (-2.7591,-1.2570)
	-- (-2.5050,-1.3966)
	-- (-2.2567,-1.5239)
	-- (-2.0137,-1.6380)
	-- (-1.7759,-1.7383)
	-- (-1.5432,-1.8244)
	-- (-1.3156,-1.8959)
	-- (-1.0933,-1.9530)
	-- (-0.8766,-1.9955)
	-- (-0.6658,-2.0237)
	-- (-0.4613,-2.0378)
	-- (-0.2636,-2.0383)
	-- (-0.0730,-2.0254)
	-- (0.1098,-1.9998)
	-- (0.2845,-1.9620)
	-- (0.4505,-1.9126)
	-- (0.6074,-1.8522)
	-- (0.7546,-1.7816)
	-- (0.8919,-1.7015)
	-- (1.0187,-1.6128)
	-- (1.1348,-1.5162)
	-- (1.2398,-1.4125)
	-- (1.3336,-1.3028)
	-- (1.4158,-1.1877)
	-- (1.4863,-1.0684)
	-- (1.5451,-0.9455)
	-- (1.5922,-0.8202)
	-- (1.6275,-0.6932)
	-- (1.6512,-0.5656)
	-- (1.6634,-0.4380)
	-- (1.6644,-0.3116)
	-- (1.6543,-0.1870)
	-- (1.6336,-0.0651)
	-- (1.6027,0.0533)
	-- (1.5619,0.1675)
	-- (1.5118,0.2767)
	-- (1.4529,0.3802)
	-- (1.3859,0.4775)
	-- (1.3113,0.5680)
	-- (1.2297,0.6511)
	-- (1.1420,0.7265)
	-- (1.0489,0.7936)
	-- (0.9510,0.8522)
	-- (0.8491,0.9020)
	-- (0.7441,0.9428)
	-- (0.6367,0.9744)
	-- (0.5278,0.9967)
	-- (0.4181,1.0097)
	-- (0.3084,1.0135)
	-- (0.1995,1.0082)
	-- (0.0922,0.9939)
	-- (-0.0129,0.9709)
	-- (-0.1149,0.9395)
	-- (-0.2132,0.9000)
	-- (-0.3071,0.8528)
	-- (-0.3961,0.7984)
	-- (-0.4795,0.7372)
	-- (-0.5568,0.6699)
	-- (-0.6276,0.5970)
	-- (-0.6914,0.5191)
	-- (-0.7479,0.4368)
	-- (-0.7967,0.3510)
	-- (-0.8376,0.2621)
	-- (-0.8704,0.1711)
	-- (-0.8950,0.0786)
	-- (-0.9112,-0.0147)
	-- (-0.9192,-0.1081)
	-- (-0.9189,-0.2007)
	-- (-0.9105,-0.2918)
	-- (-0.8941,-0.3808)
	-- (-0.8701,-0.4669)
	-- (-0.8386,-0.5495)
	-- (-0.8001,-0.6278)
	-- (-0.7549,-0.7014)
	-- (-0.7036,-0.7696)
	-- (-0.6466,-0.8318)
	-- (-0.5844,-0.8877)
	-- (-0.5178,-0.9366)
	-- (-0.4473,-0.9783)
	-- (-0.3736,-1.0124)
	-- (-0.2973,-1.0385)
	-- (-0.2193,-1.0566)
	-- (-0.1402,-1.0663)
	-- (-0.0608,-1.0676)
	-- (0.0181,-1.0604)
	-- (0.0958,-1.0448)
	-- (0.1715,-1.0208)
	-- (0.2444,-0.9886)
	-- (0.3139,-0.9484)
	-- (0.3790,-0.9004)
	-- (0.4392,-0.8450)
	-- (0.4937,-0.7826)
	-- (0.5418,-0.7135)
	-- (0.5830,-0.6384)
	-- (0.6166,-0.5577)
	-- (0.6420,-0.4720)
	-- (0.6589,-0.3820)
	-- (0.6666,-0.2883)
	-- (0.6648,-0.1917)
	-- (0.6532,-0.0929)
	-- (0.6313,0.0074)
	-- (0.5990,0.1082)
	-- (0.5561,0.2089)
	-- (0.5025,0.3086)
	-- (0.4380,0.4065)
	-- (0.3626,0.5016)
	-- (0.2764,0.5933)
	-- (0.1795,0.6806)
	-- (0.0721,0.7629)
	-- (-0.0458,0.8391)
	-- (-0.1737,0.9088)
	-- (-0.3115,0.9709)
	-- (-0.4588,1.0250)
	-- (-0.6152,1.0703)
	-- (-0.7803,1.1062)
	-- (-0.9536,1.1322)
	-- (-1.1349,1.1478)
	-- (-1.3235,1.1526)
	-- (-1.5190,1.1461)
	-- (-1.7210,1.1282)
	-- (-1.9291,1.0986)
	-- (-2.1428,1.0574)
	-- (-2.3617,1.0045)
	-- (-2.5855,0.9401)
	-- (-2.8141,0.8646)
	-- (-3.0470,0.7784)
	-- (-3.2843,0.6824)
	-- (-3.5258,0.5775)
	-- (-3.7715,0.4649)
	-- (-4.0214,0.3463)
	-- (-4.2756,0.2237)
	-- (-4.5339,0.0998)
	-- (-4.7961,-0.0225)
	-- (-5.0616,-0.1394)
	-- (-5.3291,-0.2464)
	-- (-5.5966,-0.3386)
	-- (-5.8607,-0.4105)
	-- (-6.1162,-0.4565)
	-- (-6.3560,-0.4717)
	-- (-6.5709,-0.4526)
	-- (-6.7497,-0.3982)
	-- (-6.8802,-0.3113)
	-- (-6.9513,-0.1994)
	-- (-6.9544,-0.0741)
	-- (-6.8857,0.0500)
	-- (-6.7476,0.1577)
	-- (-6.5478,0.2355)
	-- (-6.2986,0.2745)
	-- (-6.0135,0.2704)
	-- (-5.7061,0.2240)
	-- (-5.3876,0.1398)
	-- (-5.0667,0.0240)
	-- (-4.7492,-0.1157)
	-- (-4.4386,-0.2722)
	-- (-4.1368,-0.4387)
	-- (-3.8443,-0.6098)
	-- (-3.5608,-0.7805)
	-- (-3.2857,-0.9473)
	-- (-3.0184,-1.1071)
	-- (-2.7580,-1.2576)
	-- (-2.5040,-1.3971)
	-- (-2.2556,-1.5244)
	-- (-2.0127,-1.6385)
	-- (-1.7749,-1.7387)
	-- (-1.5422,-1.8247)
	-- (-1.3146,-1.8962)
	-- (-1.0924,-1.9532)
	-- (-0.8757,-1.9956)
	-- (-0.6650,-2.0238)
	-- (-0.4605,-2.0379)
	-- (-0.2628,-2.0382)
	-- (-0.0722,-2.0254)
	-- (0.1106,-1.9997)
	-- (0.2852,-1.9618)
	-- (0.4512,-1.9123)
	-- (0.6080,-1.8519)
	-- (0.7552,-1.7813)
	-- (0.8924,-1.7012)
	-- (1.0192,-1.6124)
	-- (1.1353,-1.5157)
	-- (1.2403,-1.4121)
	-- (1.3339,-1.3023)
	-- (1.4161,-1.1872)
	-- (1.4866,-1.0679)
	-- (1.5454,-0.9450)
	-- (1.5924,-0.8197)
	-- (1.6276,-0.6927)
	-- (1.6513,-0.5650)
	-- (1.6634,-0.4375)
	-- (1.6643,-0.3110)
	-- (1.6542,-0.1864)
	-- (1.6335,-0.0646)
	-- (1.6025,0.0538)
	-- (1.5617,0.1679)
	-- (1.5116,0.2771)
	-- (1.4527,0.3806)
	-- (1.3856,0.4779)
	-- (1.3109,0.5684)
	-- (1.2294,0.6515)
	-- (1.1417,0.7268)
	-- (1.0485,0.7939)
	-- (0.9505,0.8525)
	-- (0.8487,0.9022)
	-- (0.7437,0.9429)
	-- (0.6363,0.9745)
	-- (0.5273,0.9968)
	-- (0.4176,1.0098)
	-- (0.3079,1.0135)
	-- (0.1990,1.0082)
	-- (0.0917,0.9939)
	-- (-0.0133,0.9708)
	-- (-0.1153,0.9394)
	-- (-0.2136,0.8998)
	-- (-0.3075,0.8526)
	-- (-0.3964,0.7981)
	-- (-0.4798,0.7370)
	-- (-0.5571,0.6696)
	-- (-0.6279,0.5967)
	-- (-0.6917,0.5187)
	-- (-0.7481,0.4365)
	-- (-0.7969,0.3506)
	-- (-0.8378,0.2618)
	-- (-0.8705,0.1707)
	-- (-0.8951,0.0782)
	-- (-0.9113,-0.0151)
	-- (-0.9192,-0.1085)
	-- (-0.9189,-0.2010)
	-- (-0.9104,-0.2922)
	-- (-0.8941,-0.3811)
	-- (-0.8699,-0.4672)
	-- (-0.8384,-0.5498)
	-- (-0.7999,-0.6281)
	-- (-0.7547,-0.7017)
	-- (-0.7033,-0.7699)
	-- (-0.6463,-0.8321)
	-- (-0.5842,-0.8879)
	-- (-0.5175,-0.9368)
	-- (-0.4470,-0.9785)
	-- (-0.3732,-1.0125)
	-- (-0.2970,-1.0386)
	-- (-0.2190,-1.0566)
	-- (-0.1399,-1.0663)
	-- (-0.0605,-1.0676)
	-- (0.0185,-1.0604)
	-- (0.0962,-1.0447)
	-- (0.1718,-1.0207)
	-- (0.2447,-0.9885)
	-- (0.3141,-0.9482)
	-- (0.3793,-0.9002)
	-- (0.4394,-0.8448)
	-- (0.4939,-0.7823)
	-- (0.5420,-0.7132)
	-- (0.5831,-0.6381)
	-- (0.6167,-0.5573)
	-- (0.6421,-0.4716)
	-- (0.6589,-0.3816)
	-- (0.6666,-0.2879)
	-- (0.6648,-0.1913)
	-- (0.6531,-0.0924)
	-- (0.6312,0.0078)
	-- (0.5989,0.1087)
	-- (0.5559,0.2094)
	-- (0.5022,0.3090)
	-- (0.4377,0.4069)
	-- (0.3623,0.5020)
	-- (0.2761,0.5937)
	-- (0.1791,0.6810)
	-- (0.0716,0.7632)
	-- (-0.0463,0.8395)
	-- (-0.1743,0.9090)
	-- (-0.3121,0.9712)
	-- (-0.4594,1.0252)
	-- (-0.6158,1.0705)
	-- (-0.7810,1.1064)
	-- (-0.9544,1.1323)
	-- (-1.1356,1.1479)
	-- (-1.3243,1.1526)
	-- (-1.5199,1.1461)
	-- (-1.7219,1.1281)
	-- (-1.9300,1.0985)
	-- (-2.1437,1.0572)
	-- (-2.3626,1.0042)
	-- (-2.5865,0.9398)
	-- (-2.8150,0.8642)
	-- (-3.0480,0.7780)
	-- (-3.2853,0.6820)
	-- (-3.5268,0.5770)
	-- (-3.7725,0.4644)
	-- (-4.0225,0.3458)
	-- (-4.2766,0.2232)
	-- (-4.5350,0.0992)
	-- (-4.7972,-0.0230)
	-- (-5.0627,-0.1399)
	-- (-5.3303,-0.2468)
	-- (-5.5978,-0.3389)
	-- (-5.8618,-0.4107)
	-- (-6.1172,-0.4566)
	-- (-6.3570,-0.4717)
	-- (-6.5717,-0.4524)
	-- (-6.7503,-0.3979)
	-- (-6.8807,-0.3109)
	-- (-6.9515,-0.1989)
	-- (-6.9542,-0.0736)
	-- (-6.8853,0.0505)
	-- (-6.7468,0.1581)
	-- (-6.5469,0.2358)
	-- (-6.2974,0.2746)
	-- (-6.0123,0.2703)
	-- (-5.7048,0.2238)
	-- (-5.3863,0.1393)
	-- (-5.0653,0.0235)
	-- (-4.7479,-0.1163)
	-- (-4.4374,-0.2728)
	-- (-4.1356,-0.4395)
	-- (-3.8431,-0.6105)
	-- (-3.5596,-0.7813)
	-- (-3.2846,-0.9480)
	-- (-3.0173,-1.1077)
	-- (-2.7570,-1.2582)
	-- (-2.5029,-1.3977)
	-- (-2.2546,-1.5249)
	-- (-2.0117,-1.6389)
	-- (-1.7739,-1.7391)
	-- (-1.5412,-1.8250)
	-- (-1.3137,-1.8965)
	-- (-1.0915,-1.9534)
	-- (-0.8748,-1.9958)
	-- (-0.6641,-2.0239)
	-- (-0.4597,-2.0379)
	-- (-0.2620,-2.0382)
	-- (-0.0715,-2.0253)
	-- (0.1113,-1.9996)
	-- (0.2859,-1.9616)
	-- (0.4519,-1.9121)
	-- (0.6087,-1.8516)
	-- (0.7558,-1.7810)
	-- (0.8930,-1.7008)
	-- (1.0198,-1.6120)
	-- (1.1357,-1.5153)
	-- (1.2407,-1.4116)
	-- (1.3343,-1.3018)
	-- (1.4164,-1.1867)
	-- (1.4869,-1.0673)
	-- (1.5456,-0.9445)
	-- (1.5925,-0.8191)
	-- (1.6278,-0.6922)
	-- (1.6513,-0.5645)
	-- (1.6635,-0.4370)
	-- (1.6643,-0.3105)
	-- (1.6542,-0.1859)
	-- (1.6334,-0.0641)
	-- (1.6024,0.0543)
	-- (1.5615,0.1684)
	-- (1.5114,0.2776)
	-- (1.4524,0.3811)
	-- (1.3853,0.4783)
	-- (1.3106,0.5687)
	-- (1.2290,0.6518)
	-- (1.1413,0.7271)
	-- (1.0481,0.7942)
	-- (0.9501,0.8527)
	-- (0.8482,0.9024)
	-- (0.7432,0.9431)
	-- (0.6358,0.9746)
	-- (0.5269,0.9968)
	-- (0.4172,1.0098)
	-- (0.3075,1.0135)
	-- (0.1986,1.0081)
	-- (0.0913,0.9938)
	-- (-0.0137,0.9707)
	-- (-0.1157,0.9392)
	-- (-0.2140,0.8996)
	-- (-0.3079,0.8524)
	-- (-0.3968,0.7979)
	-- (-0.4802,0.7367)
	-- (-0.5575,0.6693)
	-- (-0.6282,0.5963)
	-- (-0.6919,0.5184)
	-- (-0.7483,0.4361)
	-- (-0.7971,0.3502)
	-- (-0.8379,0.2614)
	-- (-0.8707,0.1703)
	-- (-0.8952,0.0778)
	-- (-0.9113,-0.0155)
	-- (-0.9192,-0.1088)
	-- (-0.9189,-0.2014)
	-- (-0.9104,-0.2926)
	-- (-0.8940,-0.3815)
	-- (-0.8698,-0.4676)
	-- (-0.8383,-0.5501)
	-- (-0.7997,-0.6285)
	-- (-0.7545,-0.7020)
	-- (-0.7031,-0.7701)
	-- (-0.6460,-0.8323)
	-- (-0.5839,-0.8881)
	-- (-0.5172,-0.9370)
	-- (-0.4467,-0.9786)
	-- (-0.3729,-1.0126)
	-- (-0.2967,-1.0387)
	-- (-0.2186,-1.0567)
	-- (-0.1395,-1.0663)
	-- (-0.0601,-1.0676)
	-- (0.0188,-1.0603)
	-- (0.0965,-1.0446)
	-- (0.1721,-1.0206)
	-- (0.2450,-0.9883)
	-- (0.3144,-0.9480)
	-- (0.3795,-0.9000)
	-- (0.4396,-0.8445)
	-- (0.4941,-0.7820)
	-- (0.5422,-0.7129)
	-- (0.5833,-0.6377)
	-- (0.6168,-0.5570)
	-- (0.6422,-0.4712)
	-- (0.6590,-0.3812)
	-- (0.6666,-0.2875)
	-- (0.6648,-0.1909)
	-- (0.6530,-0.0920)
	-- (0.6311,0.0082)
	-- (0.5987,0.1091)
	-- (0.5557,0.2098)
	-- (0.5020,0.3094)
	-- (0.4374,0.4073)
	-- (0.3619,0.5024)
	-- (0.2757,0.5941)
	-- (0.1787,0.6814)
	-- (0.0711,0.7635)
	-- (-0.0468,0.8398)
	-- (-0.1748,0.9093)
	-- (-0.3127,0.9714)
	-- (-0.4600,1.0254)
	-- (-0.6165,1.0707)
	-- (-0.7817,1.1065)
	-- (-0.9551,1.1324)
	-- (-1.1364,1.1479)
	-- (-1.3251,1.1526)
	-- (-1.5207,1.1460)
	-- (-1.7228,1.1280)
	-- (-1.9309,1.0983)
	-- (-2.1446,1.0570)
	-- (-2.3635,1.0040)
	-- (-2.5875,0.9395)
	-- (-2.8160,0.8639)
	-- (-3.0490,0.7777)
	-- (-3.2863,0.6815)
	-- (-3.5278,0.5765)
	-- (-3.7736,0.4639)
	-- (-4.0235,0.3453)
	-- (-4.2777,0.2227)
	-- (-4.5361,0.0987)
	-- (-4.7983,-0.0235)
	-- (-5.0638,-0.1403)
	-- (-5.3314,-0.2473)
	-- (-5.5989,-0.3393)
	-- (-5.8629,-0.4110)
	-- (-6.1183,-0.4568)
	-- (-6.3579,-0.4717)
	-- (-6.5726,-0.4523)
	-- (-6.7510,-0.3976)
	-- (-6.8811,-0.3105)
	-- (-6.9516,-0.1984)
	-- (-6.9541,-0.0730)
	-- (-6.8848,0.0510)
	-- (-6.7461,0.1585)
	-- (-6.5459,0.2360)
	-- (-6.2963,0.2746)
	-- (-6.0110,0.2702)
	-- (-5.7035,0.2235)
}
\def\scaleD{0.2157}
\def\bildkurveE{
	(0.2699,-0.2928)
	-- (0.3222,-0.2348)
	-- (0.3688,-0.1744)
	-- (0.4101,-0.1119)
	-- (0.4463,-0.0476)
	-- (0.4776,0.0180)
	-- (0.5045,0.0848)
	-- (0.5269,0.1526)
	-- (0.5453,0.2211)
	-- (0.5599,0.2903)
	-- (0.5708,0.3600)
	-- (0.5782,0.4302)
	-- (0.5825,0.5007)
	-- (0.5837,0.5715)
	-- (0.5822,0.6426)
	-- (0.5781,0.7139)
	-- (0.5716,0.7856)
	-- (0.5632,0.8575)
	-- (0.5529,0.9297)
	-- (0.5411,1.0024)
	-- (0.5282,1.0754)
	-- (0.5146,1.1490)
	-- (0.5006,1.2231)
	-- (0.4869,1.2979)
	-- (0.4739,1.3733)
	-- (0.4624,1.4493)
	-- (0.4531,1.5259)
	-- (0.4471,1.6028)
	-- (0.4452,1.6797)
	-- (0.4487,1.7561)
	-- (0.4589,1.8310)
	-- (0.4768,1.9031)
	-- (0.5038,1.9707)
	-- (0.5406,2.0314)
	-- (0.5875,2.0821)
	-- (0.6439,2.1196)
	-- (0.7080,2.1403)
	-- (0.7766,2.1410)
	-- (0.8451,2.1195)
	-- (0.9079,2.0752)
	-- (0.9596,2.0099)
	-- (0.9952,1.9271)
	-- (1.0117,1.8318)
	-- (1.0078,1.7299)
	-- (0.9846,1.6265)
	-- (0.9444,1.5262)
	-- (0.8904,1.4318)
	-- (0.8260,1.3450)
	-- (0.7544,1.2665)
	-- (0.6785,1.1961)
	-- (0.6004,1.1332)
	-- (0.5218,1.0768)
	-- (0.4440,1.0260)
	-- (0.3679,0.9798)
	-- (0.2942,0.9373)
	-- (0.2231,0.8977)
	-- (0.1549,0.8604)
	-- (0.0898,0.8247)
	-- (0.0278,0.7902)
	-- (-0.0310,0.7565)
	-- (-0.0867,0.7232)
	-- (-0.1393,0.6902)
	-- (-0.1888,0.6572)
	-- (-0.2354,0.6241)
	-- (-0.2789,0.5908)
	-- (-0.3196,0.5572)
	-- (-0.3574,0.5233)
	-- (-0.3924,0.4889)
	-- (-0.4246,0.4542)
	-- (-0.4541,0.4191)
	-- (-0.4809,0.3836)
	-- (-0.5050,0.3479)
	-- (-0.5264,0.3118)
	-- (-0.5453,0.2755)
	-- (-0.5615,0.2391)
	-- (-0.5752,0.2025)
	-- (-0.5863,0.1659)
	-- (-0.5949,0.1294)
	-- (-0.6010,0.0929)
	-- (-0.6047,0.0567)
	-- (-0.6059,0.0208)
	-- (-0.6046,-0.0148)
	-- (-0.6010,-0.0499)
	-- (-0.5949,-0.0845)
	-- (-0.5866,-0.1183)
	-- (-0.5759,-0.1515)
	-- (-0.5629,-0.1837)
	-- (-0.5477,-0.2151)
	-- (-0.5303,-0.2453)
	-- (-0.5107,-0.2744)
	-- (-0.4890,-0.3023)
	-- (-0.4653,-0.3287)
	-- (-0.4395,-0.3537)
	-- (-0.4117,-0.3771)
	-- (-0.3820,-0.3987)
	-- (-0.3504,-0.4185)
	-- (-0.3170,-0.4364)
	-- (-0.2820,-0.4522)
	-- (-0.2452,-0.4658)
	-- (-0.2069,-0.4771)
	-- (-0.1670,-0.4860)
	-- (-0.1258,-0.4922)
	-- (-0.0832,-0.4957)
	-- (-0.0395,-0.4964)
	-- (0.0054,-0.4939)
	-- (0.0512,-0.4883)
	-- (0.0979,-0.4793)
	-- (0.1452,-0.4668)
	-- (0.1931,-0.4505)
	-- (0.2413,-0.4303)
	-- (0.2896,-0.4060)
	-- (0.3378,-0.3772)
	-- (0.3856,-0.3438)
	-- (0.4327,-0.3056)
	-- (0.4787,-0.2622)
	-- (0.5232,-0.2134)
	-- (0.5658,-0.1588)
	-- (0.6057,-0.0982)
	-- (0.6425,-0.0313)
	-- (0.6752,0.0422)
	-- (0.7029,0.1225)
	-- (0.7244,0.2099)
	-- (0.7384,0.3045)
	-- (0.7434,0.4061)
	-- (0.7375,0.5146)
	-- (0.7185,0.6292)
	-- (0.6841,0.7488)
	-- (0.6316,0.8717)
	-- (0.5584,0.9953)
	-- (0.4619,1.1158)
	-- (0.3399,1.2281)
	-- (0.1915,1.3260)
	-- (0.0174,1.4019)
	-- (-0.1790,1.4473)
	-- (-0.3913,1.4542)
	-- (-0.6096,1.4164)
	-- (-0.8214,1.3308)
	-- (-1.0127,1.1992)
	-- (-1.1707,1.0283)
	-- (-1.2857,0.8295)
	-- (-1.3531,0.6163)
	-- (-1.3729,0.4022)
	-- (-1.3499,0.1991)
	-- (-1.2913,0.0155)
	-- (-1.2054,-0.1434)
	-- (-1.1005,-0.2755)
	-- (-0.9837,-0.3809)
	-- (-0.8610,-0.4609)
	-- (-0.7368,-0.5180)
	-- (-0.6144,-0.5548)
	-- (-0.4962,-0.5742)
	-- (-0.3835,-0.5785)
	-- (-0.2774,-0.5701)
	-- (-0.1782,-0.5511)
	-- (-0.0862,-0.5230)
	-- (-0.0014,-0.4875)
	-- (0.0764,-0.4458)
	-- (0.1474,-0.3988)
	-- (0.2119,-0.3475)
	-- (0.2701,-0.2925)
	-- (0.3224,-0.2346)
	-- (0.3690,-0.1741)
	-- (0.4103,-0.1116)
	-- (0.4464,-0.0474)
	-- (0.4778,0.0183)
	-- (0.5046,0.0851)
	-- (0.5270,0.1529)
	-- (0.5454,0.2214)
	-- (0.5599,0.2906)
	-- (0.5708,0.3603)
	-- (0.5782,0.4305)
	-- (0.5825,0.5010)
	-- (0.5837,0.5718)
	-- (0.5822,0.6429)
	-- (0.5780,0.7142)
	-- (0.5716,0.7859)
	-- (0.5631,0.8578)
	-- (0.5528,0.9300)
	-- (0.5411,1.0027)
	-- (0.5282,1.0757)
	-- (0.5145,1.1493)
	-- (0.5006,1.2234)
	-- (0.4868,1.2982)
	-- (0.4738,1.3736)
	-- (0.4623,1.4496)
	-- (0.4531,1.5262)
	-- (0.4470,1.6031)
	-- (0.4452,1.6800)
	-- (0.4488,1.7564)
	-- (0.4589,1.8313)
	-- (0.4769,1.9034)
	-- (0.5039,1.9710)
	-- (0.5408,2.0316)
	-- (0.5877,2.0823)
	-- (0.6442,2.1198)
	-- (0.7083,2.1404)
	-- (0.7769,2.1410)
	-- (0.8453,2.1193)
	-- (0.9082,2.0750)
	-- (0.9598,2.0096)
	-- (0.9953,1.9267)
	-- (1.0117,1.8314)
	-- (1.0078,1.7294)
	-- (0.9845,1.6261)
	-- (0.9442,1.5258)
	-- (0.8901,1.4314)
	-- (0.8257,1.3447)
	-- (0.7541,1.2662)
	-- (0.6782,1.1959)
	-- (0.6000,1.1330)
	-- (0.5215,1.0766)
	-- (0.4437,1.0258)
	-- (0.3676,0.9796)
	-- (0.2939,0.9371)
	-- (0.2228,0.8976)
	-- (0.1546,0.8602)
	-- (0.0895,0.8245)
	-- (0.0276,0.7900)
	-- (-0.0313,0.7563)
	-- (-0.0869,0.7231)
	-- (-0.1395,0.6901)
	-- (-0.1890,0.6571)
	-- (-0.2355,0.6240)
	-- (-0.2791,0.5907)
	-- (-0.3198,0.5571)
	-- (-0.3576,0.5231)
	-- (-0.3926,0.4888)
	-- (-0.4248,0.4541)
	-- (-0.4542,0.4190)
	-- (-0.4810,0.3835)
	-- (-0.5051,0.3477)
	-- (-0.5265,0.3117)
	-- (-0.5454,0.2754)
	-- (-0.5616,0.2389)
	-- (-0.5753,0.2024)
	-- (-0.5864,0.1658)
	-- (-0.5950,0.1292)
	-- (-0.6011,0.0928)
	-- (-0.6047,0.0566)
	-- (-0.6059,0.0206)
	-- (-0.6046,-0.0150)
	-- (-0.6009,-0.0501)
	-- (-0.5949,-0.0846)
	-- (-0.5865,-0.1185)
	-- (-0.5758,-0.1516)
	-- (-0.5629,-0.1839)
	-- (-0.5477,-0.2152)
	-- (-0.5302,-0.2455)
	-- (-0.5107,-0.2746)
	-- (-0.4889,-0.3024)
	-- (-0.4652,-0.3288)
	-- (-0.4393,-0.3538)
	-- (-0.4116,-0.3772)
	-- (-0.3818,-0.3988)
	-- (-0.3503,-0.4186)
	-- (-0.3169,-0.4365)
	-- (-0.2818,-0.4523)
	-- (-0.2450,-0.4659)
	-- (-0.2067,-0.4772)
	-- (-0.1669,-0.4860)
	-- (-0.1256,-0.4922)
	-- (-0.0830,-0.4957)
	-- (-0.0393,-0.4963)
	-- (0.0056,-0.4939)
	-- (0.0514,-0.4883)
	-- (0.0981,-0.4793)
	-- (0.1454,-0.4668)
	-- (0.1933,-0.4505)
	-- (0.2415,-0.4302)
	-- (0.2898,-0.4058)
	-- (0.3380,-0.3771)
	-- (0.3858,-0.3437)
	-- (0.4329,-0.3054)
	-- (0.4789,-0.2620)
	-- (0.5234,-0.2131)
	-- (0.5659,-0.1586)
	-- (0.6059,-0.0980)
	-- (0.6426,-0.0310)
	-- (0.6753,0.0425)
	-- (0.7030,0.1229)
	-- (0.7245,0.2103)
	-- (0.7385,0.3049)
	-- (0.7434,0.4066)
	-- (0.7374,0.5150)
	-- (0.7184,0.6297)
	-- (0.6839,0.7493)
	-- (0.6314,0.8722)
	-- (0.5581,0.9958)
	-- (0.4614,1.1163)
	-- (0.3393,1.2286)
	-- (0.1908,1.3264)
	-- (0.0166,1.4021)
	-- (-0.1799,1.4474)
	-- (-0.3922,1.4542)
	-- (-0.6105,1.4161)
	-- (-0.8222,1.3303)
	-- (-1.0134,1.1985)
	-- (-1.1712,1.0276)
	-- (-1.2861,0.8286)
	-- (-1.3532,0.6153)
	-- (-1.3729,0.4013)
	-- (-1.3497,0.1983)
	-- (-1.2910,0.0148)
	-- (-1.2050,-0.1441)
	-- (-1.1000,-0.2760)
	-- (-0.9832,-0.3812)
	-- (-0.8604,-0.4612)
	-- (-0.7363,-0.5182)
	-- (-0.6139,-0.5550)
	-- (-0.4957,-0.5742)
	-- (-0.3831,-0.5785)
	-- (-0.2769,-0.5701)
	-- (-0.1778,-0.5510)
	-- (-0.0858,-0.5229)
	-- (-0.0010,-0.4874)
	-- (0.0767,-0.4456)
	-- (0.1477,-0.3986)
	-- (0.2122,-0.3472)
	-- (0.2704,-0.2923)
	-- (0.3226,-0.2343)
	-- (0.3692,-0.1739)
	-- (0.4104,-0.1113)
	-- (0.4466,-0.0471)
	-- (0.4779,0.0186)
	-- (0.5047,0.0854)
	-- (0.5271,0.1532)
	-- (0.5455,0.2217)
	-- (0.5600,0.2909)
	-- (0.5708,0.3606)
	-- (0.5783,0.4308)
	-- (0.5825,0.5013)
	-- (0.5837,0.5721)
	-- (0.5821,0.6432)
	-- (0.5780,0.7145)
	-- (0.5716,0.7862)
	-- (0.5631,0.8581)
	-- (0.5528,0.9303)
	-- (0.5410,1.0030)
	-- (0.5281,1.0761)
	-- (0.5145,1.1496)
	-- (0.5005,1.2238)
	-- (0.4867,1.2985)
	-- (0.4738,1.3739)
	-- (0.4623,1.4499)
	-- (0.4531,1.5265)
	-- (0.4470,1.6034)
	-- (0.4452,1.6804)
	-- (0.4488,1.7567)
	-- (0.4590,1.8316)
	-- (0.4770,1.9037)
	-- (0.5041,1.9713)
	-- (0.5409,2.0318)
	-- (0.5879,2.0825)
	-- (0.6444,2.1199)
	-- (0.7086,2.1404)
	-- (0.7772,2.1409)
	-- (0.8456,2.1192)
	-- (0.9084,2.0748)
	-- (0.9600,2.0093)
	-- (0.9954,1.9263)
	-- (1.0117,1.8310)
	-- (1.0077,1.7290)
	-- (0.9843,1.6257)
	-- (0.9440,1.5254)
	-- (0.8899,1.4310)
	-- (0.8254,1.3443)
	-- (0.7538,1.2659)
	-- (0.6778,1.1956)
	-- (0.5997,1.1327)
	-- (0.5211,1.0764)
	-- (0.4434,1.0256)
	-- (0.3673,0.9794)
	-- (0.2935,0.9370)
	-- (0.2225,0.8974)
	-- (0.1543,0.8601)
	-- (0.0892,0.8244)
	-- (0.0273,0.7899)
	-- (-0.0315,0.7562)
	-- (-0.0872,0.7229)
	-- (-0.1397,0.6899)
	-- (-0.1892,0.6570)
	-- (-0.2357,0.6239)
	-- (-0.2793,0.5905)
	-- (-0.3199,0.5569)
	-- (-0.3577,0.5230)
	-- (-0.3927,0.4886)
	-- (-0.4249,0.4539)
	-- (-0.4543,0.4188)
	-- (-0.4811,0.3833)
	-- (-0.5052,0.3476)
	-- (-0.5266,0.3115)
	-- (-0.5454,0.2752)
	-- (-0.5617,0.2388)
	-- (-0.5753,0.2022)
	-- (-0.5864,0.1656)
	-- (-0.5950,0.1291)
	-- (-0.6011,0.0926)
	-- (-0.6047,0.0564)
	-- (-0.6059,0.0205)
	-- (-0.6046,-0.0151)
	-- (-0.6009,-0.0502)
	-- (-0.5949,-0.0847)
	-- (-0.5865,-0.1186)
	-- (-0.5758,-0.1517)
	-- (-0.5628,-0.1840)
	-- (-0.5476,-0.2153)
	-- (-0.5302,-0.2456)
	-- (-0.5106,-0.2747)
	-- (-0.4889,-0.3025)
	-- (-0.4651,-0.3289)
	-- (-0.4392,-0.3539)
	-- (-0.4114,-0.3772)
	-- (-0.3817,-0.3989)
	-- (-0.3501,-0.4187)
	-- (-0.3168,-0.4366)
	-- (-0.2816,-0.4524)
	-- (-0.2449,-0.4660)
	-- (-0.2065,-0.4772)
	-- (-0.1667,-0.4860)
	-- (-0.1254,-0.4923)
	-- (-0.0829,-0.4957)
	-- (-0.0391,-0.4963)
	-- (0.0058,-0.4939)
	-- (0.0516,-0.4883)
	-- (0.0983,-0.4793)
	-- (0.1456,-0.4667)
	-- (0.1935,-0.4504)
	-- (0.2417,-0.4301)
	-- (0.2900,-0.4057)
	-- (0.3382,-0.3769)
	-- (0.3860,-0.3435)
	-- (0.4331,-0.3052)
	-- (0.4791,-0.2618)
	-- (0.5236,-0.2129)
	-- (0.5661,-0.1583)
	-- (0.6061,-0.0977)
	-- (0.6428,-0.0307)
	-- (0.6754,0.0428)
	-- (0.7031,0.1232)
	-- (0.7245,0.2107)
	-- (0.7385,0.3053)
	-- (0.7434,0.4070)
	-- (0.7374,0.5155)
	-- (0.7183,0.6302)
	-- (0.6837,0.7498)
	-- (0.6311,0.8728)
	-- (0.5577,0.9963)
	-- (0.4609,1.1168)
	-- (0.3387,1.2290)
	-- (0.1901,1.3268)
	-- (0.0159,1.4024)
	-- (-0.1807,1.4475)
	-- (-0.3931,1.4541)
	-- (-0.6115,1.4159)
	-- (-0.8231,1.3299)
	-- (-1.0142,1.1979)
	-- (-1.1718,1.0268)
	-- (-1.2865,0.8278)
	-- (-1.3534,0.6144)
	-- (-1.3729,0.4005)
	-- (-1.3496,0.1975)
	-- (-1.2907,0.0141)
	-- (-1.2046,-0.1447)
	-- (-1.0995,-0.2765)
	-- (-0.9827,-0.3816)
	-- (-0.8599,-0.4615)
	-- (-0.7357,-0.5184)
	-- (-0.6134,-0.5551)
	-- (-0.4952,-0.5743)
	-- (-0.3826,-0.5785)
	-- (-0.2765,-0.5700)
	-- (-0.1774,-0.5509)
	-- (-0.0855,-0.5228)
	-- (-0.0007,-0.4872)
	-- (0.0771,-0.4454)
	-- (0.1480,-0.3984)
	-- (0.2124,-0.3470)
	-- (0.2706,-0.2921)
	-- (0.3228,-0.2341)
	-- (0.3694,-0.1736)
	-- (0.4106,-0.1111)
	-- (0.4467,-0.0468)
	-- (0.4780,0.0188)
	-- (0.5048,0.0857)
	-- (0.5272,0.1534)
	-- (0.5455,0.2220)
	-- (0.5600,0.2912)
	-- (0.5709,0.3609)
	-- (0.5783,0.4311)
	-- (0.5825,0.5016)
	-- (0.5837,0.5724)
	-- (0.5821,0.6435)
	-- (0.5780,0.7148)
	-- (0.5716,0.7865)
	-- (0.5630,0.8584)
	-- (0.5527,0.9306)
	-- (0.5410,1.0033)
	-- (0.5281,1.0764)
	-- (0.5144,1.1499)
	-- (0.5004,1.2241)
	-- (0.4867,1.2988)
	-- (0.4737,1.3742)
	-- (0.4622,1.4502)
	-- (0.4530,1.5268)
	-- (0.4470,1.6038)
	-- (0.4452,1.6807)
	-- (0.4488,1.7570)
	-- (0.4590,1.8319)
	-- (0.4771,1.9040)
	-- (0.5042,1.9715)
	-- (0.5411,2.0321)
	-- (0.5882,2.0827)
	-- (0.6447,2.1200)
	-- (0.7089,2.1405)
	-- (0.7775,2.1409)
	-- (0.8459,2.1191)
	-- (0.9087,2.0745)
	-- (0.9602,2.0090)
	-- (0.9956,1.9259)
	-- (1.0117,1.8306)
	-- (1.0076,1.7286)
	-- (0.9842,1.6252)
	-- (0.9438,1.5250)
	-- (0.8896,1.4306)
	-- (0.8251,1.3440)
	-- (0.7535,1.2656)
	-- (0.6775,1.1953)
	-- (0.5994,1.1325)
	-- (0.5208,1.0762)
	-- (0.4431,1.0254)
	-- (0.3670,0.9793)
	-- (0.2932,0.9368)
	-- (0.2222,0.8972)
	-- (0.1540,0.8599)
	-- (0.0890,0.8242)
	-- (0.0270,0.7897)
	-- (-0.0317,0.7560)
	-- (-0.0874,0.7228)
	-- (-0.1399,0.6898)
	-- (-0.1894,0.6568)
	-- (-0.2359,0.6237)
	-- (-0.2795,0.5904)
	-- (-0.3201,0.5568)
	-- (-0.3579,0.5228)
	-- (-0.3928,0.4885)
	-- (-0.4250,0.4538)
	-- (-0.4545,0.4187)
	-- (-0.4812,0.3832)
	-- (-0.5053,0.3474)
	-- (-0.5267,0.3114)
	-- (-0.5455,0.2751)
	-- (-0.5617,0.2386)
	-- (-0.5754,0.2021)
	-- (-0.5865,0.1655)
	-- (-0.5950,0.1289)
	-- (-0.6011,0.0925)
	-- (-0.6047,0.0562)
	-- (-0.6059,0.0203)
	-- (-0.6046,-0.0153)
	-- (-0.6009,-0.0504)
	-- (-0.5948,-0.0849)
	-- (-0.5864,-0.1188)
	-- (-0.5757,-0.1519)
	-- (-0.5628,-0.1841)
	-- (-0.5475,-0.2155)
	-- (-0.5301,-0.2457)
	-- (-0.5105,-0.2748)
	-- (-0.4888,-0.3026)
	-- (-0.4650,-0.3290)
	-- (-0.4391,-0.3540)
	-- (-0.4113,-0.3773)
	-- (-0.3816,-0.3990)
	-- (-0.3500,-0.4188)
	-- (-0.3166,-0.4366)
	-- (-0.2815,-0.4524)
	-- (-0.2447,-0.4660)
	-- (-0.2064,-0.4773)
	-- (-0.1665,-0.4861)
	-- (-0.1253,-0.4923)
	-- (-0.0827,-0.4958)
	-- (-0.0389,-0.4963)
	-- (0.0060,-0.4939)
	-- (0.0518,-0.4882)
	-- (0.0985,-0.4792)
	-- (0.1458,-0.4666)
	-- (0.1937,-0.4503)
	-- (0.2419,-0.4300)
	-- (0.2902,-0.4056)
	-- (0.3384,-0.3768)
	-- (0.3862,-0.3434)
	-- (0.4333,-0.3051)
	-- (0.4793,-0.2616)
	-- (0.5238,-0.2127)
	-- (0.5663,-0.1581)
	-- (0.6062,-0.0974)
	-- (0.6429,-0.0305)
	-- (0.6756,0.0431)
	-- (0.7032,0.1236)
	-- (0.7246,0.2111)
	-- (0.7386,0.3057)
	-- (0.7434,0.4075)
	-- (0.7373,0.5160)
	-- (0.7182,0.6306)
	-- (0.6836,0.7503)
	-- (0.6309,0.8733)
	-- (0.5573,0.9968)
	-- (0.4605,1.1172)
	-- (0.3382,1.2295)
	-- (0.1894,1.3272)
	-- (0.0151,1.4026)
	-- (-0.1816,1.4476)
	-- (-0.3940,1.4541)
	-- (-0.6124,1.4156)
	-- (-0.8239,1.3294)
	-- (-1.0149,1.1972)
	-- (-1.1724,1.0260)
	-- (-1.2869,0.8269)
	-- (-1.3536,0.6135)
	-- (-1.3729,0.3996)
	-- (-1.3494,0.1967)
	-- (-1.2903,0.0133)
	-- (-1.2042,-0.1453)
	-- (-1.0991,-0.2770)
	-- (-0.9822,-0.3820)
	-- (-0.8594,-0.4617)
	-- (-0.7352,-0.5186)
	-- (-0.6129,-0.5552)
	-- (-0.4947,-0.5743)
	-- (-0.3821,-0.5785)
	-- (-0.2761,-0.5699)
	-- (-0.1770,-0.5508)
	-- (-0.0851,-0.5226)
	-- (-0.0004,-0.4870)
	-- (0.0774,-0.4452)
	-- (0.1483,-0.3982)
	-- (0.2127,-0.3468)
	-- (0.2708,-0.2918)
	-- (0.3230,-0.2338)
	-- (0.3696,-0.1733)
	-- (0.4107,-0.1108)
	-- (0.4468,-0.0465)
	-- (0.4781,0.0191)
	-- (0.5049,0.0860)
	-- (0.5273,0.1537)
	-- (0.5456,0.2223)
	-- (0.5601,0.2915)
	-- (0.5709,0.3612)
	-- (0.5783,0.4313)
	-- (0.5825,0.5019)
	-- (0.5837,0.5727)
	-- (0.5821,0.6438)
	-- (0.5780,0.7151)
	-- (0.5715,0.7868)
	-- (0.5630,0.8587)
	-- (0.5527,0.9310)
	-- (0.5409,1.0036)
	-- (0.5280,1.0767)
	-- (0.5143,1.1502)
	-- (0.5004,1.2244)
	-- (0.4866,1.2991)
	-- (0.4737,1.3745)
	-- (0.4622,1.4506)
	-- (0.4530,1.5271)
	-- (0.4470,1.6041)
	-- (0.4452,1.6810)
	-- (0.4488,1.7574)
	-- (0.4591,1.8322)
	-- (0.4772,1.9043)
	-- (0.5043,1.9718)
	-- (0.5413,2.0323)
	-- (0.5884,2.0829)
	-- (0.6449,2.1201)
	-- (0.7091,2.1405)
	-- (0.7778,2.1408)
	-- (0.8462,2.1189)
	-- (0.9089,2.0743)
	-- (0.9603,2.0086)
	-- (0.9957,1.9255)
	-- (1.0118,1.8301)
	-- (1.0076,1.7281)
	-- (0.9841,1.6248)
	-- (0.9436,1.5245)
	-- (0.8894,1.4303)
	-- (0.8248,1.3436)
	-- (0.7532,1.2653)
	-- (0.6772,1.1950)
	-- (0.5991,1.1322)
	-- (0.5205,1.0759)
	-- (0.4427,1.0252)
	-- (0.3667,0.9791)
	-- (0.2929,0.9366)
	-- (0.2219,0.8971)
	-- (0.1538,0.8598)
	-- (0.0887,0.8241)
	-- (0.0268,0.7896)
	-- (-0.0320,0.7559)
	-- (-0.0876,0.7227)
	-- (-0.1402,0.6897)
	-- (-0.1896,0.6567)
	-- (-0.2361,0.6236)
	-- (-0.2796,0.5903)
	-- (-0.3203,0.5567)
	-- (-0.3580,0.5227)
	-- (-0.3930,0.4883)
	-- (-0.4251,0.4536)
	-- (-0.4546,0.4185)
	-- (-0.4813,0.3830)
	-- (-0.5054,0.3473)
	-- (-0.5268,0.3112)
	-- (-0.5456,0.2749)
	-- (-0.5618,0.2385)
	-- (-0.5754,0.2019)
	-- (-0.5865,0.1653)
	-- (-0.5951,0.1288)
	-- (-0.6011,0.0923)
	-- (-0.6047,0.0561)
	-- (-0.6059,0.0202)
	-- (-0.6046,-0.0154)
	-- (-0.6009,-0.0505)
	-- (-0.5948,-0.0850)
	-- (-0.5864,-0.1189)
	-- (-0.5757,-0.1520)
	-- (-0.5627,-0.1843)
	-- (-0.5475,-0.2156)
	-- (-0.5300,-0.2458)
	-- (-0.5104,-0.2749)
	-- (-0.4887,-0.3027)
	-- (-0.4648,-0.3292)
	-- (-0.4390,-0.3541)
	-- (-0.4112,-0.3774)
	-- (-0.3815,-0.3991)
	-- (-0.3499,-0.4189)
	-- (-0.3165,-0.4367)
	-- (-0.2813,-0.4525)
	-- (-0.2446,-0.4661)
	-- (-0.2062,-0.4773)
	-- (-0.1663,-0.4861)
	-- (-0.1251,-0.4923)
	-- (-0.0825,-0.4958)
	-- (-0.0387,-0.4963)
	-- (0.0062,-0.4939)
	-- (0.0520,-0.4882)
	-- (0.0987,-0.4792)
	-- (0.1460,-0.4666)
	-- (0.1939,-0.4502)
	-- (0.2421,-0.4300)
	-- (0.2904,-0.4055)
	-- (0.3386,-0.3767)
	-- (0.3864,-0.3432)
	-- (0.4335,-0.3049)
	-- (0.4795,-0.2614)
	-- (0.5240,-0.2125)
	-- (0.5665,-0.1578)
	-- (0.6064,-0.0972)
	-- (0.6431,-0.0302)
	-- (0.6757,0.0435)
	-- (0.7033,0.1239)
	-- (0.7247,0.2115)
	-- (0.7386,0.3061)
	-- (0.7434,0.4079)
	-- (0.7373,0.5164)
	-- (0.7181,0.6311)
	-- (0.6834,0.7509)
	-- (0.6306,0.8738)
	-- (0.5570,0.9974)
	-- (0.4600,1.1177)
	-- (0.3376,1.2299)
	-- (0.1888,1.3275)
	-- (0.0143,1.4029)
	-- (-0.1825,1.4477)
	-- (-0.3950,1.4540)
	-- (-0.6133,1.4154)
	-- (-0.8248,1.3290)
	-- (-1.0156,1.1966)
	-- (-1.1730,1.0252)
	-- (-1.2873,0.8260)
	-- (-1.3538,0.6126)
	-- (-1.3729,0.3987)
	-- (-1.3492,0.1959)
	-- (-1.2900,0.0126)
	-- (-1.2038,-0.1459)
	-- (-1.0986,-0.2775)
	-- (-0.9817,-0.3824)
	-- (-0.8589,-0.4620)
	-- (-0.7347,-0.5188)
	-- (-0.6124,-0.5553)
	-- (-0.4942,-0.5743)
	-- (-0.3817,-0.5784)
	-- (-0.2756,-0.5699)
	-- (-0.1766,-0.5507)
	-- (-0.0847,-0.5225)
	-- (-0.0000,-0.4869)
	-- (0.0777,-0.4450)
	-- (0.1486,-0.3980)
	-- (0.2129,-0.3466)
	-- (0.2710,-0.2916)
	-- (0.3232,-0.2336)
	-- (0.3697,-0.1731)
	-- (0.4109,-0.1105)
	-- (0.4470,-0.0463)
	-- (0.4783,0.0194)
	-- (0.5050,0.0862)
	-- (0.5274,0.1540)
	-- (0.5457,0.2226)
	-- (0.5601,0.2918)
	-- (0.5710,0.3615)
	-- (0.5783,0.4316)
	-- (0.5825,0.5021)
	-- (0.5837,0.5730)
	-- (0.5821,0.6441)
	-- (0.5780,0.7154)
	-- (0.5715,0.7871)
	-- (0.5630,0.8590)
	-- (0.5527,0.9313)
	-- (0.5409,1.0039)
	-- (0.5280,1.0770)
	-- (0.5143,1.1506)
	-- (0.5003,1.2247)
	-- (0.4866,1.2995)
	-- (0.4736,1.3749)
	-- (0.4621,1.4509)
	-- (0.4530,1.5275)
	-- (0.4470,1.6044)
	-- (0.4452,1.6813)
	-- (0.4489,1.7577)
	-- (0.4591,1.8325)
	-- (0.4773,1.9046)
	-- (0.5045,1.9721)
	-- (0.5415,2.0325)
	-- (0.5886,2.0831)
	-- (0.6452,2.1203)
	-- (0.7094,2.1406)
	-- (0.7781,2.1408)
	-- (0.8465,2.1188)
	-- (0.9092,2.0741)
	-- (0.9605,2.0083)
	-- (0.9958,1.9252)
	-- (1.0118,1.8297)
	-- (1.0075,1.7277)
	-- (0.9839,1.6244)
	-- (0.9434,1.5241)
	-- (0.8891,1.4299)
	-- (0.8245,1.3433)
	-- (0.7529,1.2650)
	-- (0.6769,1.1947)
	-- (0.5987,1.1320)
	-- (0.5202,1.0757)
	-- (0.4424,1.0250)
	-- (0.3664,0.9789)
	-- (0.2926,0.9365)
	-- (0.2216,0.8969)
	-- (0.1535,0.8596)
	-- (0.0884,0.8240)
	-- (0.0265,0.7895)
	-- (-0.0322,0.7558)
	-- (-0.0878,0.7225)
	-- (-0.1404,0.6895)
	-- (-0.1898,0.6565)
	-- (-0.2363,0.6234)
	-- (-0.2798,0.5901)
	-- (-0.3204,0.5565)
	-- (-0.3582,0.5225)
	-- (-0.3931,0.4882)
	-- (-0.4253,0.4535)
	-- (-0.4547,0.4184)
	-- (-0.4814,0.3829)
	-- (-0.5055,0.3471)
	-- (-0.5269,0.3111)
	-- (-0.5457,0.2748)
	-- (-0.5618,0.2383)
	-- (-0.5755,0.2017)
	-- (-0.5865,0.1652)
	-- (-0.5951,0.1286)
	-- (-0.6011,0.0922)
	-- (-0.6047,0.0559)
	-- (-0.6059,0.0200)
	-- (-0.6046,-0.0156)
	-- (-0.6009,-0.0506)
	-- (-0.5948,-0.0852)
	-- (-0.5864,-0.1190)
	-- (-0.5756,-0.1522)
	-- (-0.5626,-0.1844)
	-- (-0.5474,-0.2157)
	-- (-0.5299,-0.2460)
	-- (-0.5103,-0.2750)
	-- (-0.4886,-0.3028)
	-- (-0.4647,-0.3293)
	-- (-0.4389,-0.3542)
	-- (-0.4111,-0.3775)
	-- (-0.3813,-0.3992)
	-- (-0.3497,-0.4189)
	-- (-0.3163,-0.4368)
	-- (-0.2812,-0.4525)
	-- (-0.2444,-0.4661)
	-- (-0.2060,-0.4774)
	-- (-0.1662,-0.4861)
	-- (-0.1249,-0.4923)
	-- (-0.0823,-0.4958)
	-- (-0.0385,-0.4963)
	-- (0.0064,-0.4939)
	-- (0.0522,-0.4882)
	-- (0.0989,-0.4791)
	-- (0.1462,-0.4665)
	-- (0.1941,-0.4502)
	-- (0.2423,-0.4299)
	-- (0.2906,-0.4054)
	-- (0.3388,-0.3766)
	-- (0.3866,-0.3431)
	-- (0.4337,-0.3047)
	-- (0.4796,-0.2612)
	-- (0.5241,-0.2123)
	-- (0.5666,-0.1576)
	-- (0.6066,-0.0969)
	-- (0.6432,-0.0299)
	-- (0.6758,0.0438)
	-- (0.7034,0.1243)
	-- (0.7248,0.2118)
	-- (0.7386,0.3066)
	-- (0.7434,0.4083)
	-- (0.7372,0.5169)
	-- (0.7180,0.6316)
	-- (0.6832,0.7514)
	-- (0.6303,0.8743)
	-- (0.5566,0.9979)
	-- (0.4595,1.1182)
	-- (0.3370,1.2304)
	-- (0.1881,1.3279)
	-- (0.0135,1.4032)
	-- (-0.1833,1.4479)
	-- (-0.3959,1.4539)
	-- (-0.6142,1.4151)
	-- (-0.8257,1.3285)
	-- (-1.0164,1.1959)
	-- (-1.1736,1.0244)
	-- (-1.2877,0.8251)
	-- (-1.3540,0.6117)
	-- (-1.3729,0.3978)
	-- (-1.3490,0.1950)
	-- (-1.2897,0.0119)
	-- (-1.2034,-0.1465)
	-- (-1.0981,-0.2780)
	-- (-0.9812,-0.3828)
	-- (-0.8584,-0.4623)
	-- (-0.7342,-0.5190)
	-- (-0.6119,-0.5554)
	-- (-0.4937,-0.5744)
	-- (-0.3812,-0.5784)
	-- (-0.2752,-0.5698)
	-- (-0.1762,-0.5506)
	-- (-0.0844,-0.5224)
	-- (0.0003,-0.4867)
	-- (0.0780,-0.4448)
	-- (0.1489,-0.3978)
	-- (0.2132,-0.3463)
	-- (0.2713,-0.2913)
	-- (0.3234,-0.2333)
	-- (0.3699,-0.1728)
	-- (0.4111,-0.1103)
	-- (0.4471,-0.0460)
	-- (0.4784,0.0197)
	-- (0.5051,0.0865)
	-- (0.5275,0.1543)
	-- (0.5457,0.2229)
	-- (0.5602,0.2921)
	-- (0.5710,0.3618)
	-- (0.5784,0.4319)
	-- (0.5825,0.5024)
	-- (0.5837,0.5733)
	-- (0.5821,0.6444)
	-- (0.5779,0.7157)
	-- (0.5715,0.7874)
	-- (0.5629,0.8593)
	-- (0.5526,0.9316)
	-- (0.5408,1.0042)
	-- (0.5279,1.0773)
	-- (0.5142,1.1509)
	-- (0.5003,1.2250)
	-- (0.4865,1.2998)
	-- (0.4736,1.3752)
	-- (0.4621,1.4512)
	-- (0.4529,1.5278)
	-- (0.4470,1.6047)
	-- (0.4452,1.6817)
	-- (0.4489,1.7580)
	-- (0.4592,1.8329)
	-- (0.4774,1.9049)
	-- (0.5046,1.9723)
	-- (0.5417,2.0328)
	-- (0.5888,2.0833)
	-- (0.6454,2.1204)
	-- (0.7097,2.1406)
	-- (0.7783,2.1407)
	-- (0.8467,2.1186)
	-- (0.9094,2.0738)
	-- (0.9607,2.0080)
	-- (0.9959,1.9248)
	-- (1.0118,1.8293)
	-- (1.0075,1.7272)
	-- (0.9838,1.6240)
	-- (0.9432,1.5237)
}
\def\scaleE{0.6574}

\def\l{0.10}
\def\punkt#1{
	\fill[color=blue!20] #1 circle[radius=\l];
	\fill[color=white] #1 circle[radius={0.5*\l}];
	\draw[color=blue] #1 circle[radius={0.5*\l}];
}
\def\kreuz#1{
	\fill[color=blue!20] #1 circle[radius=\l];
	\foreach \w in {45,135,...,360}{
		\draw[color=blue] #1 -- ++(\w:\l);
	}
}
\def\achsen{
	\draw[->] (-1.6,0) -- (1.7,0)
		coordinate[label={$\operatorname{Re}$}];
	\draw[->] (0,-1.6) -- (0,1.6)
		coordinate[label={right:$\operatorname{Im}$}];
}
\def\f{
	\draw[->] (-0.5,0) -- ++(1.2,0);
	\node at (0.1,0) [above] {$f(z)$};
}
\begin{tikzpicture}[>=latex,thick,scale=\skala]

%
% 1 Nullstelle
%
\begin{scope}[yshift=14.0cm]
	\begin{scope}[xshift=-2.8cm]
		\achsen
		\draw[color=darkred,line width=1.2pt] (0,0) circle[radius=\r];
		\punkt{\nullstelleeins}
	\end{scope}
	\f
	\begin{scope}[xshift=2.8cm]
		\achsen
		\begin{scope}[scale=\scaleA]
			\draw[color=darkred,line width=1.2pt] \bildkurveA;
		\end{scope}
		%\node at (1.6,-1.6) [above left] {$n=1$};
	\end{scope}
\end{scope}

%
% 2 Nullstellen
%
\begin{scope}[yshift=10.5cm]
	\begin{scope}[xshift=-2.8cm]
		\achsen
		\draw[color=darkred,line width=1.2pt] (0,0) circle[radius=\r];
		\punkt{\nullstelleeins}
		\punkt{\nullstellezwei}
	\end{scope}
	\f
	\begin{scope}[xshift=2.8cm]
		\achsen
		\begin{scope}[scale=\scaleB]
			\draw[color=darkred,line width=1.2pt] \bildkurveB;
		\end{scope}
		%\node at (1.6,-1.6) [above left] {$n=3$};
	\end{scope}
\end{scope}

%
% 3 Nullstellen
%
\begin{scope}[yshift=7.0cm]
	\begin{scope}[xshift=-2.8cm]
		\achsen
		\draw[color=darkred,line width=1.2pt] (0,0) circle[radius=\r];
		\punkt{\nullstelleeins}
		\punkt{\nullstellezwei}
		\punkt{\nullstelledrei}
	\end{scope}
	\f
	\begin{scope}[xshift=2.8cm]
		\achsen
		\begin{scope}[scale=\scaleC]
			\draw[color=darkred,line width=1.2pt] \bildkurveC;
		\end{scope}
		%\node at (1.6,-1.6) [above left] {$n=3$};
	\end{scope}
\end{scope}

%
% 3 Nullstelle, 1 Polstelle innerhalb
%
\begin{scope}[yshift=3.5cm]
	\begin{scope}[xshift=-2.8cm]
		\achsen
		\draw[color=darkred,line width=1.2pt] (0,0) circle[radius=\r];
		\punkt{\nullstelleeins}
		\punkt{\nullstellezwei}
		\punkt{\nullstelledrei}
		\kreuz{\poleins}
	\end{scope}
	\f
	\begin{scope}[xshift=2.8cm]
		\achsen
		\begin{scope}[scale=\scaleD]
			\draw[color=darkred,line width=1.2pt] \bildkurveD;
		\end{scope}
		%\node at (1.6,-1.6) [above left] {$n=2$};
	\end{scope}
\end{scope}

%
% 3 Nullstelle, 1 Polstelle innerhalb, 1 Polstelle ausserhalb
%
\begin{scope}
	\begin{scope}[xshift=-2.8cm]
		\achsen
		\draw[color=darkred,line width=1.2pt] (0,0) circle[radius=\r];
		\punkt{\nullstelleeins}
		\punkt{\nullstellezwei}
		\punkt{\nullstelledrei}
		\kreuz{\poleins}
		\kreuz{\polzwei}
	\end{scope}
	\f
	\begin{scope}[xshift=2.8cm]
		\achsen
		\begin{scope}[scale=\scaleE]
			\draw[color=darkred,line width=1.2pt] \bildkurveE;
		\end{scope}
		%\node at (1.6,-1.6) [above left] {$n=2$};
	\end{scope}
\end{scope}

\end{tikzpicture}
\end{document}

